\documentclass[12pt, a4paper]{article}
\usepackage[utf8]{inputenc}
\usepackage[spanish]{babel}
\usepackage{geometry}
\usepackage{hyperref}
\usepackage{longtable}
\usepackage{array}
\usepackage{xcolor}
\usepackage{enumitem}
\usepackage{booktabs}
\usepackage{multirow}
\usepackage{times}
\usepackage{fancyhdr}
\usepackage{titlesec}
\usepackage{tabularx}

\geometry{top=2.54cm, bottom=2.54cm, left=2.54cm, right=2.54cm}
\renewcommand{\rmdefault}{ptm}
\renewcommand{\sfdefault}{phv}
\renewcommand{\baselinestretch}{1.08}
\setlength{\parindent}{0.5cm}
\setlength{\parskip}{0.4cm}

\definecolor{criticolor}{RGB}{180,0,0}
\definecolor{medio}{RGB}{30,100,180}

\pagestyle{fancy}
\fancyhf{}
\rhead{\small SRS Qurakuna APP v5.2}
\lhead{\small ISO/IEC/IEEE 29148:2018}
\rfoot{\small P\'agina \thepage}
\lfoot{\small UAC -- Fondos Concursables 2025-I}

\hypersetup{
  colorlinks=true,
  linkcolor=blue,
  urlcolor=blue,
  pdftitle={SRS Qurakuna APP v5.2}
}

\newcommand{\criticomark}{\textcolor{criticolor}{\textbf{[CR\'ITICO]}}}
\newcommand{\inferido}{\textcolor{medio}{\textit{[INFERIDO]}}}

\title{
  \textbf{Especificaci\'on de Requerimientos de Software (SRS)\\
  para el Sistema \textit{Qurakuna APP}}\\[0.5em]
  \large Versi\'on 5.2 -- Basado en ISO/IEC/IEEE 29148:2018
}
\author{
  Universidad Andina del Cusco (UAC)\\
  Proyecto de Investigaci\'on -- Fondos Concursables 2025-I
}
\date{Mayo 2026}

\begin{document}

\maketitle
\thispagestyle{empty}

\begin{center}
\small
\begin{tabular}{|c|c|p{6.5cm}|c|}
\hline
\textbf{Versi\'on} & \textbf{Fecha} & \textbf{Cambios principales} & \textbf{Autor} \\
\hline
1.0 & 20/03/2026 & Versi\'on inicial (m\'odulos 1--3) & Equipo UAC \\
\hline
2.0 & 23/03/2026 & M\'odulos 4--9, matriz de trazabilidad, RBAC, Actor 4, RNF cuantificados & Equipo UAC \\
\hline
3.0 & 24/03/2026 & RF-017 eliminado, RF-016 reubicado, RNF-013 elevado, RNF-018 por especie & Equipo UAC \\
\hline
4.0 & 28/03/2026 & Modelo de datos normalizado, token\_local, merge RF-026, umbral Laplaciano remoto & Equipo UAC \\
\hline
5.0 & 14/04/2026 & CR-01 a CR-06: 11 atributos Ochoa, fallback GPS, purga fotos, zoom offline, OTA resumable, RNF-020 bater\'ia & Equipo UAC \\
\hline
5.1 & 22/04/2026 & Sistema h\'ibrido app+web; flujo asistido RF-049/050/051; mapa de calor; 16 especies; nueva especie & Equipo UAC \\
\hline
5.2 & 26/04/2026 & Roles clarificados (4 actores: Guardaparque, Validador, Bot\'anico, Administrador); stack web Node.js+Express confirmado; mapa de observaciones simple en app y web reemplaza mapa de calor; RF-046 web integrado; RNF interoperabilidad y paleta corporativa; Ap\'endice A actualizado & Equipo UAC \\
\hline
5.2-rev1 & 09/05/2026 & Hardware: RAM como piso tecnológico; RF-003: cifrado de token con Keystore; nuevos RF-059 (Edge AI <2s), RNF-023 (cuantización Int8, <150 MB RAM), RF-060 (cola de sincronización asíncrona); arquitectura Offline-First explícita; matriz trazabilidad actualizada & Equipo UAC \\
\hline
\end{tabular}
\end{center}

\newpage
\tableofcontents
\newpage

%=============================================================
\section{Introducci\'on}
%=============================================================

\subsection{Prop\'osito del documento}

El presente documento define los requerimientos de software del sistema \textbf{Qurakuna APP} conforme al est\'andar \textbf{ISO/IEC/IEEE 29148:2018} y las categor\'ias de calidad de \textbf{ISO/IEC 25010:2011}. El sistema comprende dos componentes interoperables:

\begin{itemize}
  \item \textbf{Aplicaci\'on m\'ovil Android (campo):} para guardaparques del SERNANP, con funcionamiento offline-first.
  \item \textbf{Plataforma web (gesti\'on):} para Validadores, Bot\'anicos y el Administrador, con acceso desde navegador.
\end{itemize}

Ambos componentes comparten un \'unico \textbf{API Backend (Node.js + Express + PostgreSQL)} que act\'ua como n\'ucleo conector y fuente de verdad centralizada.

Este documento est\'a dirigido a:
\begin{itemize}
  \item \textbf{Equipo de desarrollo:} Estudiantes de Ingenier\'ia de Sistemas de la UAC.
  \item \textbf{Investigadores y asesores:} Responsable de investigaci\'on y asesora Ing. Gabriela Zu\~niga Rojas.
  \item \textbf{Personal del SERNANP:} Guardaparques, Validadores, Bot\'anico y Administrador.
  \item \textbf{Direcci\'on de Gesti\'on de la Investigaci\'on UAC:} Entidad evaluadora y financiadora.
\end{itemize}

\subsection{Alcance del sistema}

\textbf{Componente 1 -- App m\'ovil Android (campo):}
\begin{itemize}
  \item Identificaci\'on offline de las 16 especies de orqu\'ideas del SHMP mediante CNN embebida.
  \item Flujo asistido cuando la confianza del modelo es insuficiente: recomendaci\'on de nueva foto, preguntas guiadas, reporte de especie no identificada.
  \item Cat\'alogo de 16 especies con fichas completas (11 atributos Ochoa + atributos visuales).
  \item Registro de observaciones con informaci\'on biol\'ogica de campo y georreferenciaci\'on.
  \item Mapa de observaciones b\'asico con puntos georreferenciados y zonas de inter\'es del SHMP.
  \item Sincronizaci\'on diferida con el API Backend al recuperar conexi\'on.
\end{itemize}

\textbf{Componente 2 -- Plataforma web (gesti\'on):}
\begin{itemize}
  \item Dashboard de validaci\'on para el Validador (Especialista en RRNN).
  \item Edici\'on del cat\'alogo y revisi\'on de reportes de especie desconocida para el Bot\'anico.
  \item Gesti\'on de cuentas de usuario, actualizaci\'on OTA del modelo y administraci\'on general para el Administrador.
  \item Mapa de observaciones interactivo con capas de zonas de inter\'es y puestos de control.
  \item Reportes exportables en PDF.
\end{itemize}

\textbf{Componente 3 -- API Backend (Node.js + Express + PostgreSQL):}
\begin{itemize}
  \item N\'ucleo conector que recibe peticiones de la app m\'ovil y de la plataforma web.
  \item Expone \textit{endpoints} REST con estructuras JSON compatibles con la app m\'ovil.
  \item Gesti\'on de autenticaci\'on JWT, autorizaci\'on RBAC y cifrado de contrase\~nas con Bcrypt.
\end{itemize}

\textbf{Fuera del alcance:} iOS, funcionalidades colaborativas para turistas, integraci\'on en tiempo real con sistemas externos del SERNANP, entrenamiento continuo autom\'atico desde la app.

\subsection{Definiciones, acr\'onimos y abreviaciones}

\begin{longtable}{|p{3.2cm}|p{11.3cm}|}
\hline
\textbf{T\'ermino} & \textbf{Definici\'on} \\
\hline
\endhead
\textbf{Accuracy} & Proporci\'on de identificaciones correctas sobre el total. M\'etrica objetivo 85\% (RNF-018). \\
\hline
API & Application Programming Interface. Interfaz REST JSON del backend Node.js. \\
\hline
CNN & Red Neuronal Convolucional. Motor de identificaci\'on de las 16 especies en TF Lite. \\
\hline
CRISP-DM & Cross-Industry Standard Process for Data Mining. Marco metodol\'ogico del proyecto. \\
\hline
\textbf{Degradaci\'on Elegante} & El sistema ofrece la mejor alternativa disponible (ej. GPS $>$10m) sin bloquear al usuario. \\
\hline
GPS & Global Positioning System. \\
\hline
Ground truth & Datos validados por el Bot\'anico para entrenamiento y evaluaci\'on del modelo. \\
\hline
IUCN & International Union for Conservation of Nature. \\
\hline
JWT & JSON Web Token. Mecanismo de autenticaci\'on entre clientes (app, web) y API. \\
\hline
\textbf{Labelo} & P\'etalo modificado de las orqu\'ideas; elemento clave de identificaci\'on morfol\'ogica. \\
\hline
OTA & Over-The-Air. Actualizaci\'on del modelo CNN sin reinstalar la APK. \\
\hline
RBAC & Role-Based Access Control. Control de acceso por rol. \\
\hline
RF & Requerimiento Funcional. \\
\hline
RNF & Requerimiento No Funcional. \\
\hline
\textbf{Resumabilidad} & Capacidad de una descarga de reanudar desde el punto de interrupci\'on. \\
\hline
SHMP & Santuario Hist\'orico de Machu Picchu. \\
\hline
SERNANP & Servicio Nacional de \'Areas Naturales Protegidas por el Estado del Per\'u. \\
\hline
SRS & Software Requirements Specification. \\
\hline
\textbf{Tipo de Florecimiento} & Caracter\'istica de la inflorescencia: apical, lateral, basal, racemoso, solitario. \\
\hline
TF Lite & TensorFlow Lite. Framework Google para inferencia de modelos ML en m\'oviles. \\
\hline
Top-N & N resultados con mayor probabilidad de clasificaci\'on. \\
\hline
UAC & Universidad Andina del Cusco. \\
\hline
\textbf{Zoom Level} & Nivel de detalle del mapa digital. ZL 15/16 = sector/trocha del SHMP. \\
\hline
\textbf{Ochoa (11 atributos)} & 11 campos de ficha de especie de la gu\'ia de Julio Gustavo Ochoa Estrada. \\
\hline
\textbf{Precision} (ML) & Proporci\'on de identificaciones de una especie que son correctas. \\
\hline
\textbf{Recall} (ML) & Proporci\'on de ejemplos reales de una especie que el modelo identifica. \\
\hline
\textbf{F1-score} & Media arm\'onica de precision y recall. \\
\hline
\textbf{Nueva Especie (posible)} & Planta no clasificable en las 16 especies tras m\'ultiples intentos y flujo asistido. \\
\hline
\end{longtable}

\subsection{Convenciones del documento}

\begin{itemize}
  \item \textbf{Deber\'a / Debe:} SHALL -- obligatorio.
  \item \textbf{Deber\'ia:} SHOULD -- recomendado con justificaci\'on t\'ecnica.
  \item \textbf{Podr\'a:} MAY -- opcional deseable.
\end{itemize}

\subsection{Referencias}

\begin{enumerate}
  \item Qurakuna ID Requisitos.docx -- Equipo InnovaTec UAC. v1.0, feb. 2026.
  \item ANEXO 2 -- Estructura del Proyecto de Investigaci\'on 2025-I. VRIN-UAC, ago. 2025.
  \item ANEXO 3 -- Cronograma del Proyecto 2025-I. VRIN-UAC, ago. 2025.
  \item Plantas -- Aplicaci\'on de Marco de Trabajo CRISP-DM. Equipo UAC, ene. 2026.
  \item Anexo 04 -- Estructura de Informe Parcial F3-PSC-F-04. DGI-UAC.
  \item ISO/IEC/IEEE 29148:2018 -- Requirements Engineering.
  \item ISO/IEC 25010:2011 -- Software Quality Requirements and Evaluation.
  \item SERNANP (2023). Reporte anual del SHMP. Ministerio del Ambiente del Per\'u.
  \item Ochoa Estrada, J.G. -- Gu\'ia de orqu\'ideas del SHMP (11 atributos est\'andar de ficha).
  \item Documento de Arquitectura y Requisitos Detallados (React, Node.js y App M\'ovil) -- Equipo UAC Web, abr. 2026.
\end{enumerate}

\subsection{Visi\'on general del documento}

Secci\'on 1: introducci\'on y alcance. Secci\'on 2: actores, caracter\'isticas del sistema y restricciones. 
Secci\'on 3: RF-001 a RF-060 en 10 m\'odulos. Secci\'on 4: RNF-001 a RNF-023. Secci\'on 5: ap\'endices (modelo de datos, casos de uso, trazabilidad, cronograma, glosario, cambios).

%=============================================================
\section{Descripci\'on general}
%=============================================================

\subsection{Perspectiva del producto}

El sistema \textbf{Qurakuna APP} est\'a compuesto por tres capas interoperables:

\begin{enumerate}
  \item \textbf{App m\'ovil Android} (Kotlin + TF Lite): opera offline-first; sincroniza con el API al recuperar conexi\'on.
  \item \textbf{Plataforma web} (React.js 18 + Tailwind CSS 3): interfaz de gesti\'on accesible desde navegador con conexi\'on.
  \item \textbf{API Backend} (Node.js 20 LTS + Express + PostgreSQL 15): n\'ucleo conector unificado; expone endpoints REST JSON compatibles con la app m\'ovil y la plataforma web.
\end{enumerate}

La app m\'ovil se comunica con la API mediante Axios/Fetch API. La plataforma web tambi\'en consume el mismo API. El dise\~no de las respuestas del API est\'a dictado por lo que la app m\'ovil necesita recibir, garantizando compatibilidad estricta (Ref. 10).

Se adopta una \textbf{Arquitectura Offline-First con Inferencia en el Borde (Edge)}. El backend (Node.js/PostgreSQL) act\'ua como repositorio central de conocimiento, mientras que la l\'ogica de predicci\'on reside en el cliente m\'ovil para asegurar la operatividad en las zonas de sombra de red del Santuario Hist\'orico de Machu Picchu.

\subsection{Caracter\'isticas del sistema}

\begin{longtable}{|p{0.8cm}|p{3.5cm}|p{7cm}|p{1.8cm}|}
\hline
\textbf{ID} & \textbf{Feature} & \textbf{Descripci\'on} & \textbf{Prior.} \\
\hline
\endhead
F1 & Autenticaci\'on y roles & RBAC para 4 roles; JWT; sesi\'on persistente offline (app) & Alta \\
\hline
F2 & Captura de im\'agenes & C\'amara, validaci\'on de calidad, preprocesamiento & Alta \\
\hline
F3 & Identificaci\'on IA & CNN offline, 16 especies, top-3 con endemismo e IUCN & Alta \\
\hline
F3b & Flujo asistido & Recomendaci\'on adaptativa, preguntas guiadas, reporte & Alta \\
\hline
F4 & Cat\'alogo de especies & 16 fichas (11 atributos Ochoa + visuales), offline & Alta \\
\hline
F5 & Registro de observaciones & GPS degradado, datos de campo, purga post-sync & Alta \\
\hline
F6 & Validaci\'on (web) & Dashboard: aprobar/corregir/rechazar observaciones & Alta \\
\hline
F7 & Mapa de observaciones & Puntos georreferenciados + zonas de inter\'es SHMP (app y web) & Media \\
\hline
F8 & Reportes & PDF exportable, m\'etricas, encuesta usabilidad & Media \\
\hline
F9 & Administraci\'on (web) & Gesti\'on de cuentas, OTA modelo, cat\'alogo, API Backend & Alta \\
\hline
F10 & Plataforma web & Dashboard integrado para Validador, Bot\'anico y Admin & Alta \\
\hline
\end{longtable}

\subsection{Actores del sistema}

El sistema tiene \textbf{cuatro actores}, con plataforma y responsabilidades claramente particionadas:

\subsubsection{Actor 1: Guardaparque (App m\'ovil -- campo)}
\begin{longtable}{|p{3.5cm}|p{11cm}|}
\hline
Tipo & Guardaparque del SERNANP \\
\hline
Plataforma & \textbf{Aplicaci\'on m\'ovil Android exclusivamente} \\
\hline
Formaci\'on & Gesti\'on tur\'istica e infraestructura. Formaci\'on m\'inima en biolog\'ia. \\
\hline
Responsabilidades & Patrullajes, captura fotogr\'afica, identificaci\'on y registro de orqu\'ideas en campo. \\
\hline
Frecuencia & Diaria, m\'ultiples sesiones por patrullaje. \\
\hline
Necesidades clave & Identificaci\'on $<$5 min, offline 100\%, flujo m\'ax. 3 pasos, saber exactamente qu\'e est\'a viendo (especie, si es end\'emica o amenazada). \\
\hline
Permisos & Captura, identificaci\'on + flujo asistido, cat\'alogo, mapa, historial propio, registro de observaciones. \\
\hline
\end{longtable}

\subsubsection{Actor 2: Validador -- Especialista en Recursos Naturales (Plataforma web)}
\begin{longtable}{|p{3.5cm}|p{11cm}|}
\hline
Tipo & Especialista en biodiversidad del SERNANP \\
\hline
Plataforma & \textbf{Plataforma web exclusivamente} (requiere conexi\'on) \\
\hline
Formaci\'on & Profesional con experiencia en gesti\'on de biodiversidad y conservaci\'on. \\
\hline
Responsabilidades & Validar identificaciones de los guardaparques, supervisar calidad de datos, analizar observaciones, retroalimentar al guardaparque. \\
\hline
Frecuencia & Diaria o semanal desde estaci\'on de trabajo. \\
\hline
Necesidades clave & Dashboard de observaciones pendientes, herramientas de correcci\'on, estad\'isticas de precisi\'on, mapa de observaciones, reportes exportables. \\
\hline
Permisos & Dashboard de validaci\'on, mapa de observaciones web, estad\'isticas de precisi\'on, exportaci\'on de reportes. No puede gestionar usuarios ni actualizar el modelo. \\
\hline
\end{longtable}

\subsubsection{Actor 3: Bot\'anico Especialista (Plataforma web)}
\begin{longtable}{|p{3.5cm}|p{11cm}|}
\hline
Tipo & Bi\'ologo tax\'onomo con especializaci\'on en orqu\'ideas andinas \\
\hline
Plataforma & \textbf{Plataforma web exclusivamente} (requiere conexi\'on) \\
\hline
Formaci\'on & Bi\'ologo con experiencia en Orchidaceae del SHMP. Contratado por S/.2,000. \\
\hline
Responsabilidades & Validar y editar las 16 fichas del cat\'alogo (11 atributos + atributos visuales), etiquetar el dataset de entrenamiento (ground truth), revisar reportes de especie desconocida. \\
\hline
Frecuencia & Fases 2 y 3 del proyecto (nov. 2025 -- abr. 2026). \\
\hline
Necesidades clave & Editor de fichas de especie con todos los atributos (labelo, florecimiento, hojas), cola de reportes de especie no identificada con im\'agenes y datos de campo. \\
\hline
Permisos & Edici\'on de fichas del cat\'alogo, revisi\'on de reportes de especie desconocida. No puede gestionar usuarios ni actualizar el modelo. \\
\hline
\end{longtable}

\subsubsection{Actor 4: Administrador del Sistema (Plataforma web)}
\begin{longtable}{|p{3.5cm}|p{11cm}|}
\hline
Tipo & Personal t\'ecnico designado por SERNANP/UAC \\
\hline
Plataforma & \textbf{Plataforma web exclusivamente} (requiere conexi\'on) \\
\hline
Formaci\'on & Conocimientos en administraci\'on de sistemas inform\'aticos. \\
\hline
Responsabilidades & \textbf{Gesti\'on de cuentas de usuario} (\'unico actor con acceso a CRUD de usuarios); actualizaci\'on OTA del modelo CNN; mantenimiento del cat\'alogo; monitoreo del sistema y del API Backend. \\
\hline
Frecuencia & Mensual o ante eventos espec\'ificos (nuevo guardaparque, nuevo modelo). \\
\hline
Necesidades clave & Panel de administraci\'on web con CRUD de usuarios, control de versiones del modelo, gesti\'on del cat\'alogo de 16 especies, panel de estad\'isticas globales. \\
\hline
Permisos & Acceso completo a todos los m\'odulos de la plataforma web. \'Unico actor que puede crear, editar y desactivar cuentas de usuario de cualquier rol. \\
\hline
\end{longtable}

\subsection{Restricciones generales}

\begin{itemize}
  \item \textbf{App m\'ovil:} Android 8.0 (API 26)+; TF Lite; APK $\leq$150 MB. El sistema deberá ser compatible con dispositivos móviles Android cuyo hardware cuente con un mínimo de 2 GB de memoria RAM disponible. Se optimizará el software para garantizar fluidez en este umbral base, aprovechando capacidades superiores en dispositivos de gama media y alta.
  \item \textbf{Plataforma web:} React.js 18.x; Tailwind CSS 3.x; compatible con Chrome/Firefox/Edge (\'ultimas 2 versiones). Paleta corporativa: \texttt{\#757B47}, \texttt{\#FEFAE0} y tonos complementarios.
  \item \textbf{API Backend:} Node.js 20.x LTS; Express; PostgreSQL 15.x; JWT; Bcrypt. Estructura JSON de respuestas compatible con la app m\'ovil (Ref. 10).
  \item \textbf{De tiempo:} Ago. 2025 -- Jul. 2026. Desarrollo: abr.--may. 2026. Validaci\'on campo: jun. 2026.
  \item \textbf{Presupuestarias:} Total S/.10,991.00. Servidor S/.500 (sincronizaci\'on lotes diferidos; $\leq$5 usuarios web concurrentes).
  \item \textbf{Legales:} Autorizaci\'on SERNANP. Ley N.° 29733 de Protecci\'on de Datos Personales.
  \item \textbf{Modelo IA:} 16 especies; 500--700 im\'agenes/especie; precisi\'on depende del dataset.
  \item \textbf{Energ\'eticas:} CNN + GPS + C\'amara concurrentes limitan la bater\'ia (RNF-020).
\end{itemize}

\subsection{Suposiciones y dependencias}

\begin{itemize}
  \item Los guardaparques tendr\'an dispositivos Android compatibles del SERNANP.
  \item El Bot\'anico validar\'a las 16 fichas (ene.--feb. 2026).
  \item La plataforma web ser\'a accesible desde los puestos de control con conexi\'on a internet.
  \item El equipo de desarrollo web tiene acceso a la documentaci\'on t\'ecnica o puede realizar ingenier\'ia inversa a las peticiones de red de la app m\'ovil para replicar \textit{endpoints} (Ref. 10, secci\'on 6).
  \item La confirmaci\'on taxon\'omica de una posible nueva especie requiere an\'alisis ex-situ por el Bot\'anico.
\end{itemize}

%=============================================================
\section{Requerimientos funcionales}
%=============================================================

Seg\'un ISO/IEC/IEEE 29148:2018, cada RF incluye: ID, nombre, obligatoriedad, descripci\'on, rationale, fuente, prioridad, verificaci\'on, stakeholder y criterios de aceptaci\'on.

%-------------------------------------------
\subsection{M\'odulo 1: Autenticaci\'on y Gesti\'on de Usuarios}
%-------------------------------------------

\textit{Plataforma: App m\'ovil (Guardaparque) + Plataforma web (Validador, Bot\'anico, Administrador). El registro y la gesti\'on de cuentas son exclusivos del Administrador desde la web.}

\begin{longtable}{|p{3cm}|p{11.5cm}|}
\hline
\multicolumn{2}{|c|}{\textbf{RF-001} \criticomark} \\
\hline
Nombre & Registro de nuevo usuario (Administrador -- Web) \\
\hline
Obligatoriedad & \textbf{deber\'a} \\
\hline
Descripci\'on & El sistema deber\'a permitir al Administrador, desde la plataforma web, registrar nuevos usuarios con: nombres completos, email institucional, rol (guardaparque / validador / bot\'anico / administrador), puesto de control asignado (para guardaparques) y credenciales temporales autogeneradas. La plataforma web env\'ia los datos en formato JSON al API Backend (HTTP 201 -- Ref. 10, RF-01). \\
\hline
Rationale & El Administrador es el \'unico actor autorizado para crear cuentas. El registro nunca es autom\'atico ni lo hacen los propios usuarios. \\
\hline
Fuente & \inferido{}. Ref. 10 (RF-01, RF-03). \\
\hline
Prioridad & Alta \\
\hline
Verificaci\'on & Inspecci\'on de red (DevTools): petici\'on POST desde React al API retorna HTTP 201. \\
\hline
Stakeholder & Administrador \\
\hline
Criterios de Aceptaci\'on &
AC-001.1: Admin crea usuario con todos los campos obligatorios completados. \newline
AC-001.2: Emails duplicados rechazados con mensaje espec\'ifico. \newline
AC-001.3: Contrase\~na temporal \'unica $\geq$8 caracteres generada autom\'aticamente. \newline
AC-001.4: Usuario registrado en PostgreSQL del servidor; si es Guardaparque, disponible en app al sincronizarse. \\
\hline
\end{longtable}

\begin{longtable}{|p{3cm}|p{11.5cm}|}
\hline
\multicolumn{2}{|c|}{\textbf{RF-002} \criticomark} \\
\hline
Nombre & Inicio de sesi\'on con credenciales \\
\hline
Obligatoriedad & \textbf{deber\'a} \\
\hline
Descripci\'on & Ambas plataformas deber\'an autenticar usuarios mediante email y contrase\~na. El API Backend valida el token JWT en cada petici\'on protegida; peticiones sin token v\'alido retornan HTTP 401 (Ref. 10, RF-03). En el primer inicio con contrase\~na temporal, el sistema exige cambio obligatorio. \\
\hline
Rationale & Autenticaci\'on segura unificada para app y web. \\
\hline
Fuente & \inferido{}. Ref. 10 (RF-03). \\
\hline
Prioridad & Alta \\
\hline
Verificaci\'on & Prueba funcional en app y en navegador web; prueba de rechazo de petici\'on sin token (HTTP 401). \\
\hline
Stakeholder & Todos los actores \\
\hline
Criterios de Aceptaci\'on &
AC-002.1: Credenciales correctas en app: acceso $\leq$3 segundos. \newline
AC-002.2: Credenciales correctas en web: acceso $\leq$2 segundos. \newline
AC-002.3: Credenciales incorrectas: mensaje gen\'erico. \newline
AC-002.4: Primer inicio: fuerza cambio de contrase\~na. \newline
AC-002.5: 5 intentos fallidos: cuenta bloqueada 15 minutos. \newline
AC-002.6: Petici\'on sin JWT al API: HTTP 401. \\
\hline
\end{longtable}

\newpage
\begin{longtable}{|p{3cm}|p{11.5cm}|}
\hline
\multicolumn{2}{|c|}{\textbf{RF-003} \criticomark} \\
\hline
Nombre & Sesi\'on persistente offline (App m\'ovil -- Guardaparque) \\
\hline
Obligatoriedad & \textbf{deber\'a} \\
\hline
Descripci\'on & La app m\'ovil deber\'a mantener la sesi\'on activa sin conexi\'on hasta cierre expl\'icito o 30 d\'ias de inactividad. Si expira offline, continua operaci\'on hasta la pr\'oxima conexi\'on, momento en que fuerza reautenticaci\'on antes de sincronizar. El token de sesi\'on y cualquier credencial persistente se almacenarán de forma segura utilizando \textbf{EncryptedSharedPreferences} (o similar) respaldado por el \textbf{Android Keystore System}; no se permite almacenamiento en texto plano. \\
\hline
Rationale & Zonas del SHMP sin cobertura de red. El Guardaparque no puede autenticarse repetidamente en campo. \\
\hline
Fuente & Plantas -- CRISP-DM (Fase 1). \\
\hline
Prioridad & Alta \\
\hline
Verificaci\'on & Prueba funcional en modo avi\'on. \\
\hline
Stakeholder & Guardaparque \\
\hline
Criterios de Aceptaci\'on &
AC-003.1: App abre sin conexi\'on: acceso directo sin login. \newline
AC-003.2: ``Cerrar sesi\'on'' invalida el token local. \newline
AC-003.3: Token expirado offline: app funciona; al conectar fuerza reautenticaci\'on. \newline
AC-003.4: Campos \texttt{token\_local} y \texttt{fecha\_expiracion\_sesion} actualizados en cada inicio/cierre. \newline
AC-003.5: El \texttt{token\_local} y cualquier credencial persistente se almacenan cifrados en reposo mediante el Keystore de Android; no existen archivos de texto plano con datos de sesi\'on. \\
\hline
\end{longtable}

\begin{longtable}{|p{3cm}|p{11.5cm}|}
\hline
\multicolumn{2}{|c|}{\textbf{RF-004}} \\
\hline
Nombre & Gesti\'on de perfil de usuario \\
\hline
Obligatoriedad & \textbf{deber\'ia} \\
\hline
Descripci\'on & El sistema deber\'ia permitir a los usuarios autenticados (en app o web seg\'un su rol) visualizar y editar su perfil: cambiar contrase\~na, tel\'efono y preferencias de notificaci\'on. \\
\hline
Fuente & \inferido{} \\
\hline
Prioridad & Media \\
\hline
Criterios de Aceptaci\'on &
AC-004.1: Acceso al perfil desde men\'u principal. \newline
AC-004.2: Cambios guardados (app: sincronizan al conectar; web: inmediato). \\
\hline
\end{longtable}

\begin{longtable}{|p{3cm}|p{11.5cm}|}
\hline
\multicolumn{2}{|c|}{\textbf{RF-005} \criticomark} \\
\hline
Nombre & Control de acceso basado en roles (RBAC) \\
\hline
Obligatoriedad & \textbf{deber\'a} \\
\hline
Descripci\'on & El sistema deber\'a implementar RBAC con los siguientes permisos, con partici\'on estricta de plataformas por rol:\newline
\textbf{Guardaparque (solo app m\'ovil):} captura, identificaci\'on + flujo asistido, cat\'alogo, mapa de la app, historial propio, registro de observaciones.\newline
\textbf{Validador (solo web):} dashboard de observaciones pendientes, aprobar/corregir/rechazar, estad\'isticas de precisi\'on, mapa web, exportaci\'on de reportes.\newline
\textbf{Bot\'anico (solo web):} editor del cat\'alogo (11 atributos + atributos visuales), cola de reportes de especie desconocida. Sin acceso a gesti\'on de usuarios.\newline
\textbf{Administrador (solo web):} todos los permisos anteriores + CRUD de usuarios, actualizaci\'on OTA del modelo, gesti\'on del cat\'alogo de 16 especies, panel de estad\'isticas globales, configuraci\'on del sistema. \\
\hline
Rationale & La partici\'on de plataformas por rol resuelve el vac\'io de versiones anteriores. Los roles de gesti\'on (Validador, Bot\'anico, Admin) no necesitan app m\'ovil; el Guardaparque no necesita plataforma web. \\
\hline
Fuente & \inferido{}. Ref. 10. \\
\hline
Prioridad & Alta \\
\hline
Verificaci\'on & Prueba de permisos con sesi\'on de cada rol en app y web. \\
\hline
Stakeholder & Todos los actores \\
\hline
Criterios de Aceptaci\'on &
AC-005.1: Guardaparque no puede acceder a la plataforma web (HTTP 403). \newline
AC-005.2: Validador/Bot\'anico/Admin no pueden autenticarse en la app m\'ovil. \newline
AC-005.3: Validador accede a validaci\'on, estad\'isticas y mapa web; no puede editar el cat\'alogo. \newline
AC-005.4: Bot\'anico edita fichas y revisa reportes; no puede gestionar usuarios. \newline
AC-005.5: Administrador es el \'unico actor que crea, edita y desactiva cuentas de usuario. \\
\hline
\end{longtable}

%-------------------------------------------
\subsection{M\'odulo 2: Captura y Procesamiento de Im\'agenes}
%-------------------------------------------

\textit{Plataforma: App m\'ovil exclusivamente.}

\begin{longtable}{|p{3cm}|p{11.5cm}|}
\hline
\multicolumn{2}{|c|}{\textbf{RF-006} \criticomark} \\
\hline
Nombre & Captura desde c\'amara en tiempo real \\
\hline
Obligatoriedad & \textbf{deber\'a} \\
\hline
Descripci\'on & La app deber\'a permitir capturar fotograf\'ias con visor en tiempo real, gu\'ias de encuadre y enfoque t\'actil. C\'amara activada bajo demanda y liberada al completar la captura (RNF-020). \\
\hline
Fuente & Qurakuna ID Requisitos.docx \\
\hline
Prioridad & Alta \\
\hline
Verificaci\'on & Prueba funcional en dispositivo f\'isico. \\
\hline
Stakeholder & Guardaparque \\
\hline
Criterios de Aceptaci\'on &
AC-006.1: C\'amara abre en $\leq$2 segundos. \newline
AC-006.2: Gu\'ias de encuadre visibles. \newline
AC-006.3: Enfoque t\'actil funcional. \newline
AC-006.4: Imagen guardada temporalmente en almacenamiento local. \\
\hline
\end{longtable}

\begin{longtable}{|p{3cm}|p{11.5cm}|}
\hline
\multicolumn{2}{|c|}{\textbf{RF-007}} \\
\hline
Nombre & Carga de im\'agenes desde galer\'ia \\
\hline
Obligatoriedad & \textbf{deber\'ia} \\
\hline
Descripci\'on & La app deber\'ia permitir seleccionar una imagen existente (.jpg o .png) de la galer\'ia para identificaci\'on. \\
\hline
Fuente & Qurakuna ID Requisitos.docx \\
\hline
Prioridad & Media \\
\hline
Criterios de Aceptaci\'on &
AC-007.1: Selector de archivos del sistema se abre. \newline
AC-007.2: Solo .jpg y .png aceptados. \\
\hline
\end{longtable}

\begin{longtable}{|p{3cm}|p{11.5cm}|}
\hline
\multicolumn{2}{|c|}{\textbf{RF-008} \criticomark} \\
\hline
Nombre & Validaci\'on autom\'atica de calidad de imagen \\
\hline
Obligatoriedad & \textbf{deber\'a} \\
\hline
Descripci\'on & La app deber\'a validar autom\'aticamente: resoluci\'on m\'inima 512$\times$512 px; nitidez por varianza del Laplaciano $\geq$ \texttt{umbral\_laplaciano} (par\'ametro remoto configurable, por defecto 100); brillo promedio en [30, 220]. \\
\hline
Fuente & Qurakuna ID Requisitos.docx \\
\hline
Prioridad & Alta \\
\hline
Verificaci\'on & Prueba con conjunto de im\'agenes de calidades controladas ($\geq$20). \\
\hline
Stakeholder & Guardaparque \\
\hline
Criterios de Aceptaci\'on &
AC-008.1: $<$512$\times$512: ``Resoluci\'on insuficiente''. \newline
AC-008.2: Varianza $<$ umbral: ``Imagen borrosa, mantenga el dispositivo firme''. \newline
AC-008.3: Brillo $<$30: ``Poca luz''; brillo $>$220: ``Sobreexposici\'on''. \newline
AC-008.4: \texttt{umbral\_laplaciano} descargado al iniciar sesi\'on con conexi\'on; sin conexi\'on usa \'ultimo valor o 100. \\
\hline
\end{longtable}

\begin{longtable}{|p{3cm}|p{11.5cm}|}
\hline
\multicolumn{2}{|c|}{\textbf{RF-009} \criticomark} \\
\hline
Nombre & Preprocesamiento local de imagen \\
\hline
Obligatoriedad & \textbf{deber\'a} \\
\hline
Descripci\'on & La app deber\'a preprocesar la imagen validada: redimensionar a 224$\times$224 px con \textit{center crop} y normalizar p\'ixeles a [0, 1]. Totalmente local, sin red. \\
\hline
Fuente & Plantas -- CRISP-DM (Fase 3). \\
\hline
Prioridad & Alta \\
\hline
Criterios de Aceptaci\'on &
AC-009.1: Imagen resultante: 224$\times$224 px. \newline
AC-009.2: P\'ixeles en [0.0, 1.0] (float32). \newline
AC-009.3: Preprocesamiento en $\leq$500 ms. \newline
AC-009.4: Sin solicitudes de red. \\
\hline
\end{longtable}

\begin{longtable}{|p{3cm}|p{11.5cm}|}
\hline
\multicolumn{2}{|c|}{\textbf{RF-010}} \\
\hline
Nombre & Feedback visual de calidad al usuario \\
\hline
Obligatoriedad & \textbf{deber\'ia} \\
\hline
Descripci\'on & La app deber\'ia mostrar indicadores en tiempo real durante la captura: nivel de luz, alerta de desenfoque. Tras rechazo, mensaje con sugerencia espec\'ifica. \\
\hline
Fuente & \inferido{} \\
\hline
Prioridad & Media \\
\hline
Criterios de Aceptaci\'on &
AC-010.1: Poca luz: \'icono de alerta en visor. \newline
AC-010.2: Borrosidad: mensaje espec\'ifico en $\leq$1 segundo. \\
\hline
\end{longtable}

\begin{longtable}{|p{3cm}|p{11.5cm}|}
\hline
\multicolumn{2}{|c|}{\textbf{RF-011}} \\
\hline
Nombre & Soporte a m\'ultiples \'angulos de captura \\
\hline
Obligatoriedad & \textbf{deber\'ia} \\
\hline
Descripci\'on & La app deber\'ia recomendar capturar la orqu\'idea desde diferentes \'angulos: frontal del labelo, lateral, superior, hoja. Cada \'angulo genera una imagen independiente asociada a la observaci\'on (tabla \texttt{ImagenObservacion}, relaci\'on 1-a-N). \\
\hline
Fuente & Qurakuna ID Requisitos.docx \\
\hline
Prioridad & Media \\
\hline
Criterios de Aceptaci\'on &
AC-011.1: App muestra sugerencias de \'angulos con \'iconos ilustrativos. \\
\hline
\end{longtable}

%-------------------------------------------
\subsection{M\'odulo 3: Identificaci\'on de Especies (Motor IA)}
%-------------------------------------------

\textit{Plataforma: App m\'ovil exclusivamente. CNN offline embebida para 16 especies.}

\begin{longtable}{|p{3cm}|p{11.5cm}|}
\hline
\multicolumn{2}{|c|}{\textbf{RF-012} \criticomark} \\
\hline
Nombre & Clasificaci\'on autom\'atica mediante CNN embebida (16 especies) \\
\hline
Obligatoriedad & \textbf{deber\'a} \\
\hline
Descripci\'on & La app deber\'a ejecutar localmente un modelo CNN en formato TF Lite para clasificar la imagen preprocesada entre las \textbf{16 especies de orqu\'ideas del SHMP} validadas por el Bot\'anico. Modelo embebido en la APK; inferencia 100\% offline. \\
\hline
Fuente & Anexo 04. Plantas -- CRISP-DM (Fase 4). \\
\hline
Prioridad & Alta \\
\hline
Verificaci\'on & Prueba de integraci\'on + validaci\'on de precisi\'on con dataset de 16 clases. \\
\hline
Stakeholder & Guardaparque \\
\hline
Criterios de Aceptaci\'on &
AC-012.1: Modelo carga en $\leq$2 segundos desde almacenamiento local. \newline
AC-012.2: Produce vector de probabilidades para las 16 clases. \newline
AC-012.3: Inferencia completa en modo avi\'on. \\
\hline
\end{longtable}

\begin{longtable}{|p{3cm}|p{11.5cm}|}
\hline
\multicolumn{2}{|c|}{\textbf{RF-013} \criticomark} \\
\hline
Nombre & Mostrar top-3 con informaci\'on completa de endemismo y conservaci\'on \\
\hline
Obligatoriedad & \textbf{deber\'a} \\
\hline
Descripci\'on & La app deber\'a mostrar los 3 resultados m\'as probables con: nombre cient\'ifico, nombre com\'un (ES/EN), nombre en quechua (cuando disponible), porcentaje de confianza, barra visual y, de forma \textbf{destacada}: (a) indicador de \textbf{endemismo} (end\'emica SHMP / del Per\'u / no end\'emica) y (b) \textbf{estado IUCN} (CR / EN / VU / LC). El Guardaparque debe saber exactamente qu\'e est\'a viendo y si la especie requiere cuidado especial. \\
\hline
Rationale & Los Guardaparques necesitan conocer si la planta es end\'emica o est\'a amenazada para actuar correctamente en campo. Esta informaci\'on es cr\'itica para sus decisiones de conservaci\'on. \\
\hline
Fuente & Bot\'anico especialista. ANEXO 2. \\
\hline
Prioridad & Alta \\
\hline
Verificaci\'on & Prueba con im\'agenes de las 16 especies, incluyendo end\'emicas y amenazadas. \\
\hline
Stakeholder & Guardaparque \\
\hline
Criterios de Aceptaci\'on &
AC-013.1: 3 especies ordenadas de mayor a menor confianza. \newline
AC-013.2: Nombre cient\'ifico y porcentaje num\'erico en cada resultado. \newline
AC-013.3: Indicador de endemismo con \'icono diferenciado (SHMP / Per\'u). \newline
AC-013.4: Estado IUCN con color: rojo (CR), naranja (EN), amarillo (VU), verde (LC). \newline
AC-013.5: Informaci\'on correcta para las 16 especies seg\'un validaci\'on del Bot\'anico. \\
\hline
\end{longtable}

\begin{longtable}{|p{3cm}|p{11.5cm}|}
\hline
\multicolumn{2}{|c|}{\textbf{RF-014} \criticomark} \\
\hline
Nombre & Funcionamiento offline completo del motor IA \\
\hline
Obligatoriedad & \textbf{deber\'a} \\
\hline
Descripci\'on & La app deber\'a garantizar que captura, preprocesamiento, clasificaci\'on CNN y resultados funcionen al 100\% sin conexi\'on. Modelo y metadatos de las 16 especies embebidos. \\
\hline
Fuente & Plantas -- CRISP-DM (Fase 1). ANEXO 2. \\
\hline
Prioridad & Alta \\
\hline
Verificaci\'on & Prueba modo avi\'on con 20 identificaciones consecutivas. \\
\hline
Criterios de Aceptaci\'on &
AC-014.1: Modo avi\'on: toda la funcionalidad de identificaci\'on disponible. \newline
AC-014.2: Sin errores de red durante la identificaci\'on. \\
\hline
\end{longtable}

\begin{longtable}{|p{3cm}|p{11.5cm}|}
\hline
\multicolumn{2}{|c|}{\textbf{RF-015} \criticomark} \\
\hline
Nombre & Umbral m\'inimo de confianza configurable \\
\hline
Obligatoriedad & \textbf{deber\'a} \\
\hline
Descripci\'on & La app deber\'a implementar un umbral de confianza configurable por el Administrador (por defecto: 70\%). Si el resultado de mayor confianza es $<$ umbral, activa autom\'aticamente el flujo asistido (RF-049). \\
\hline
Fuente & \inferido{} -- CRISP-DM Fase 4. \\
\hline
Prioridad & Alta \\
\hline
Criterios de Aceptaci\'on &
AC-015.1: Confianza $<$ umbral: activa el flujo asistido RF-049. \newline
AC-015.2: Admin cambia umbral entre 50\% y 90\% desde la web; se aplica en pr\'oxima sincronizaci\'on. \\
\hline
\end{longtable}

\begin{longtable}{|p{3cm}|p{11.5cm}|}
\hline
\multicolumn{2}{|c|}{\textbf{RF-016} \textit{(Reubicado a RNF-019)}} \\
\hline
\multicolumn{2}{|p{14.5cm}|}{Soporte a estados fenol\'ogicos. Reubicado a RNF-019 (calidad del modelo IA). Ver RNF-019.} \\
\hline
\end{longtable}

\begin{longtable}{|p{3cm}|p{11.5cm}|}
\hline
\multicolumn{2}{|c|}{\textbf{RF-017} \textit{(Consolidado en RNF-001)}} \\
\hline
\multicolumn{2}{|p{14.5cm}|}{Tiempo de respuesta de identificaci\'on. Consolidado en RNF-001 para evitar duplicado. Trazabilidad $\rightarrow$ RNF-001.} \\
\hline
\end{longtable}

\begin{longtable}{|p{3cm}|p{11.5cm}|}
\hline
\multicolumn{2}{|c|}{\textbf{RF-018} \criticomark} \\
\hline
Nombre & Indicaci\'on de especies end\'emicas y en riesgo \\
\hline
Obligatoriedad & \textbf{deber\'a} \\
\hline
Descripci\'on & La app deber\'a mostrar \'icono distintivo y alerta para especies end\'emicas del Per\'u o clasificadas en IUCN como VU, EN o CR. Integrado con RF-013 en la pantalla de resultados. \\
\hline
Fuente & ANEXO 2 (IUCN 2023). Bot\'anico especialista. \\
\hline
Prioridad & Alta \\
\hline
Criterios de Aceptaci\'on &
AC-018.1: Especie CR: \'icono rojo y texto de alerta. \newline
AC-018.2: Especie end\'emica no amenazada: \'icono diferenciado. \\
\hline
\end{longtable}

% --- NUEVO RF-059 (Edge AI inmediato) ---
\begin{longtable}{|p{3cm}|p{11.5cm}|}
\hline
\multicolumn{2}{|c|}{\textbf{RF-059} \criticomark} \\
\hline
Nombre & Identificación de Especies Offline (Edge AI) \\
\hline
Obligatoriedad & \textbf{deberá} \\
\hline
Descripción & La aplicación móvil deberá realizar la clasificación de orquídeas de forma local y autónoma (Edge AI) utilizando un motor de inferencia TensorFlow Lite. La respuesta de identificación debe ser inmediata (< 2 segundos) sin requerir acceso a red en el momento de la captura. \\
\hline
Rationale & Elimina cualquier dependencia de API externa; garantiza la operación en las zonas sin cobertura del SHMP. \\
\hline
Fuente & Requisito de arquitectura Offline-First. \\
\hline
Prioridad & Alta (Crítica) \\
\hline
Verificación & Prueba en modo avión: captura → resultado en <2 s, sin tráfico de red. \\
\hline
Stakeholder & Guardaparque \\
\hline
Criterios de Aceptación & AC-059.1: Tiempo de inferencia menor a 2 segundos en dispositivo con 2 GB de RAM. \newline

AC-059.2: La identificación no genera ninguna petición HTTP durante la captura. \newline
AC-059.3: Funciona con el modelo cuantizado (Int8) embebido en la APK. \\
\hline
\end{longtable}

%-------------------------------------------
\subsection{M\'odulo 3B: Flujo Asistido de Identificaci\'on}
%-------------------------------------------

\textit{Plataforma: App m\'ovil exclusivamente. Se activa cuando la confianza CNN no alcanza el umbral (RF-015) o cuando el Guardaparque solicita ayuda.}

\begin{longtable}{|p{3cm}|p{11.5cm}|}
\hline
\multicolumn{2}{|c|}{\textbf{RF-049} \criticomark} \\
\hline
Nombre & Recomendaci\'on adaptativa de nueva fotograf\'ia \\
\hline
Obligatoriedad & \textbf{deber\'a} \\
\hline
Descripci\'on & Cuando la confianza es inferior al umbral, la app deber\'a: (a) mostrar los resultados disponibles con advertencia ``Identificaci\'on incierta''; (b) recomendar nueva fotograf\'ia con sugerencias espec\'ificas de \'angulo y condiciones: ``Enfoca el labelo frontalmente'', ``Captura tambi\'en las hojas'', ``Busca mejor iluminaci\'on''. Registra el n\'umero de intentos del ciclo. \\
\hline
Rationale & Una foto de mejor \'angulo puede superar el umbral. La sugerencia espec\'ifica orienta al Guardaparque hacia los elementos clave de identificaci\'on (labelo, hojas, tipo de florecimiento). \\
\hline
Fuente & Bot\'anico especialista. \\
\hline
Prioridad & Alta \\
\hline
Verificaci\'on & Prueba con imagen de baja confianza: recomendaci\'on en $\leq$2 segundos. \\
\hline
Stakeholder & Guardaparque \\
\hline
Criterios de Aceptaci\'on &
AC-049.1: Confianza $<$ umbral: advertencia y recomendaci\'on de nueva foto en $\leq$2 segundos. \newline
AC-049.2: Recomendaci\'on incluye al menos un \'angulo espec\'ifico (labelo, hoja o florecimiento). \newline
AC-049.3: Contador de intentos del ciclo registrado internamente. \\
\hline
\end{longtable}

\begin{longtable}{|p{3cm}|p{11.5cm}|}
\hline
\multicolumn{2}{|c|}{\textbf{RF-050} \criticomark} \\
\hline
Nombre & Identificaci\'on asistida por preguntas guiadas \\
\hline
Obligatoriedad & \textbf{deber\'a} \\
\hline
Descripci\'on & Si tras \textbf{2 o m\'as intentos fotogr\'aficos} la confianza sigue por debajo del umbral, la app deber\'a activar preguntas guiadas sobre caracter\'isticas visuales clave:\newline
\textbf{Preguntas predefinidas:}\newline
\quad -- ``\textit{?`Cu\'al es el color predominante del labelo?}'' (blanco / morado / amarillo / rosa / bicolor)\newline
\quad -- ``\textit{?`C\'omo son las hojas?}'' (lanceoladas largas / ovales cortas / en roseta basal / carnosas)\newline
\quad -- ``\textit{?`D\'onde surge la inflorescencia?}'' (escapo central / lateral / entre hojas)\newline
\quad -- ``\textit{?`D\'onde crece la planta?}'' (en el suelo / sobre rocas / sobre \'arboles)\newline
\quad -- ``\textit{?`Hay varias flores en el mismo tallo?}'' (varias flores -- racemosa / una sola flor)\newline
\textbf{Campo libre:} el Guardaparque puede ingresar observaciones adicionales (hasta 500 caracteres): color, tama\~no, condiciones, \'epoca.\newline
Respuestas almacenadas localmente y sincronizadas al servidor junto con la observaci\'on. \\
\hline
Rationale & Los atributos visuales de las orqu\'ideas (labelo, florecimiento, morfolog\'ia foliar, biotipo) son los mismos que usa el Bot\'anico en la identificaci\'on taxon\'omica. Las preguntas reproducen el proceso cognitivo del experto de forma accesible para el Guardaparque. \\
\hline
Fuente & Bot\'anico especialista. Gu\'ia Ochoa (16 especies). \\
\hline
Prioridad & Alta \\
\hline
Verificaci\'on & Prueba funcional con guardaparques simulando identificaci\'on fallida. \\
\hline
Stakeholder & Guardaparque \\
\hline
Criterios de Aceptaci\'on &
AC-050.1: Tras 2 intentos fallidos: flujo de preguntas activado autom\'aticamente. \newline
AC-050.2: Al menos 4 de las 5 preguntas con opciones visuales (\'iconos + texto). \newline
AC-050.3: Campo libre permite hasta 500 caracteres. \newline
AC-050.4: Respuestas almacenadas localmente y sincronizadas con la observaci\'on. \newline
AC-050.5: Flujo de preguntas funciona completamente offline. \\
\hline
\end{longtable}

\begin{longtable}{|p{3cm}|p{11.5cm}|}
\hline
\multicolumn{2}{|c|}{\textbf{RF-051} \criticomark} \\
\hline
Nombre & Reporte de especie desconocida / posible nueva especie \\
\hline
Obligatoriedad & \textbf{deber\'a} \\
\hline
Descripci\'on & Si tras el flujo de preguntas guiadas (RF-050) el Guardaparque no puede identificar la planta, la app deber\'a ofrecer la opci\'on de generar un \textbf{Reporte de Especie No Identificada} con:\newline
(a) Todas las fotograf\'ias capturadas (m\'ultiples \'angulos).\newline
(b) Respuestas al flujo de preguntas guiadas.\newline
(c) Coordenadas GPS (con margen real si aplica degradaci\'on elegante).\newline
(d) Fecha, hora y puesto de control.\newline
(e) Estado fenol\'ogico observado.\newline
(f) Observaciones de texto libre adicionales.\newline
(g) Vector de probabilidades del \'ultimo intento CNN.\newline
El reporte se sincroniza al servidor y aparece en la cola de revisi\'on del Bot\'anico (RF-055). Si 3 o m\'as Guardaparques distintos reportan el mismo esp\'ecimen (zona $\leq$500m, per\'iodo $\leq$30 d\'ias), el sistema genera una alerta de ``Posible Nueva Especie'' para revisi\'on prioritaria por el Bot\'anico y el Administrador. \\
\hline
Rationale & Un fallo completo de identificaci\'on tras m\'ultiples intentos es una se\~nal de alto valor cient\'ifico que debe escalarse al experto taxon\'omico. La coincidencia de m\'ultiples reportes en la misma zona aumenta la probabilidad de que sea un esp\'ecimen genuinamente no catalogado. \\
\hline
Fuente & Bot\'anico especialista. ANEXO 2 (conservaci\'on de biodiversidad). \\
\hline
Prioridad & Alta \\
\hline
Verificaci\'on & Prueba simulando m\'ultiples intentos fallidos; verificar generaci\'on y sincronizaci\'on del reporte. \\
\hline
Stakeholder & Guardaparque, Bot\'anico, Administrador \\
\hline
Criterios de Aceptaci\'on &
AC-051.1: Opci\'on ``Reportar como no identificada'' disponible tras flujo de preguntas. \newline
AC-051.2: Reporte incluye los 7 campos obligatorios. \newline
AC-051.3: Reporte almacenado localmente y sincronizado al servidor al conectar. \newline
AC-051.4: Reporte aparece en la cola de revisi\'on del Bot\'anico en la web. \newline
AC-051.5: 3+ reportes en zona/per\'iodo: alerta ``Posible Nueva Especie'' en dashboard del Bot\'anico y Administrador. \\
\hline
\end{longtable}

%-------------------------------------------
\subsection{M\'odulo 4: Cat\'alogo de Especies}
%-------------------------------------------

\textit{Plataforma: App m\'ovil (consulta offline) + Plataforma web (edici\'on, Bot\'anico/Administrador).}

\begin{longtable}{|p{3cm}|p{11.5cm}|}
\hline
\multicolumn{2}{|c|}{\textbf{RF-019} \criticomark} \\
\hline
Nombre & Ficha detallada por especie (16 especies -- 11 atributos Ochoa + atributos visuales) \\
\hline
Obligatoriedad & \textbf{deber\'a} \\
\hline
Descripci\'on & El sistema deber\'a mantener una ficha completa por cada una de las \textbf{16 especies} con los \textbf{11 atributos Ochoa}:\newline
(1) Biotipo(s) o Forma(s) de crecimiento. (2) Sistema(s) Ecol\'ogico(s). (3) Endemismo [local e* / nacional e]. (4) Estado Conservaci\'on IUCN. (5) Reporte Nuevo [Machupicchu\textsuperscript{(m)} / Per\'u\textsuperscript{(?)}]. (6) Nombre com\'un ES. (7) Descripci\'on ES. (8) Detalle flor/inflorescencia/fruto. (9) Fotograf\'ia general. (10) Nombre com\'un EN. (11) Descripci\'on EN.\newline
Adicionalmente, los siguientes \textbf{atributos visuales clave de identificaci\'on}:\newline
\quad -- \textbf{Color y morfolog\'ia del labelo:} color(es), manchas, forma (tubular/abierto/espatulado).\newline
\quad -- \textbf{Tipo de florecimiento:} apical/lateral/basal; racemoso/solitario; \'epoca t\'ipica.\newline
\quad -- \textbf{Morfolog\'ia foliar:} forma, textura, disposici\'on, longitud aproximada.\newline
\quad -- \textbf{Biotipo detallado:} ep\'ifita/terrestre/lit\'ofita.\newline
\quad -- Indicadores de presencia en Camino Inca, Jard\'in Bot\'anico Machupicchu y Jard\'in Bot\'anico de la Llaqta. \\
\hline
Rationale & Los atributos visuales del labelo, tipo de florecimiento y hojas son los elementos primarios de identificaci\'on visual de orqu\'ideas. El Guardaparque necesita saber exactamente qu\'e est\'a viendo: estos atributos alimentan las preguntas guiadas (RF-050) y permiten la comparaci\'on en campo. \\
\hline
Fuente & Bot\'anico especialista. Gu\'ia Ochoa (Ref. 9). CR-01. ANEXO 2. \\
\hline
Prioridad & Alta \\
\hline
Verificaci\'on & Revisi\'on con Bot\'anico: 11 atributos Ochoa + atributos visuales para las 16 especies. \\
\hline
Stakeholder & Guardaparque, Validador, Bot\'anico \\
\hline
Criterios de Aceptaci\'on &
AC-019.1: Ficha de cada una de las 16 especies con 11 atributos Ochoa completos. \newline
AC-019.2: Atributos visuales (labelo, florecimiento, hojas) completos para las 16 especies. \newline
AC-019.3: Campos opcionales visibles solo cuando est\'an disponibles. \newline
AC-019.4: Ficha marcada ``validada'' por el Bot\'anico. \newline
AC-019.5: BD almacena todos los atributos en la entidad \texttt{Especie} (Ap\'endice A). \\
\hline
\end{longtable}

\begin{longtable}{|p{3cm}|p{11.5cm}|}
\hline
\multicolumn{2}{|c|}{\textbf{RF-020} \criticomark} \\
\hline
Nombre & Galer\'ia de im\'agenes de referencia por especie \\
\hline
Obligatoriedad & \textbf{deber\'a} \\
\hline
Descripci\'on & Cada ficha deber\'a incluir $\geq$3 fotos de referencia por especie: frontal flor (\'enfasis en labelo), lateral e imagen de hojas. Validadas por el Bot\'anico. \\
\hline
Fuente & Qurakuna ID Requisitos.docx. CRISP-DM Fase 3. \\
\hline
Prioridad & Alta \\
\hline
Criterios de Aceptaci\'on &
AC-020.1: Cada especie tiene $\geq$3 im\'agenes (frontal-labelo, lateral, hoja). \newline
AC-020.2: Im\'agenes visibles en modo offline. \\
\hline
\end{longtable}

\begin{longtable}{|p{3cm}|p{11.5cm}|}
\hline
\multicolumn{2}{|c|}{\textbf{RF-021}} \\
\hline
Nombre & B\'usqueda y filtrado de especies \\
\hline
Obligatoriedad & \textbf{deber\'a} \\
\hline
Descripci\'on & La app deber\'a permitir buscar entre las 16 especies por nombre cient\'ifico, nombre com\'un (ES/EN) o quechua, y filtrar por: estado IUCN, zona del Santuario, estado fenol\'ogico, biotipo y sistema ecol\'ogico. \\
\hline
Fuente & \inferido{} \\
\hline
Prioridad & Alta \\
\hline
Criterios de Aceptaci\'on &
AC-021.1: B\'usqueda parcial retorna resultados en $\leq$1 segundo. \newline
AC-021.2: Filtro por IUCN funcional. \newline
AC-021.3: B\'usqueda offline. \newline
AC-021.4: Nombre exacto retorna especie correcta como primer resultado. \\
\hline
\end{longtable}

\begin{longtable}{|p{3cm}|p{11.5cm}|}
\hline
\multicolumn{2}{|c|}{\textbf{RF-022} \criticomark} \\
\hline
Nombre & Acceso completo al cat\'alogo en modo offline \\
\hline
Obligatoriedad & \textbf{deber\'a} \\
\hline
Descripci\'on & El cat\'alogo completo de las 16 especies (fichas + atributos visuales + im\'agenes de referencia) deber\'a estar disponible localmente sin conexi\'on. \\
\hline
Fuente & CRISP-DM Fase 6. ANEXO 2. \\
\hline
Prioridad & Alta \\
\hline
Criterios de Aceptaci\'on &
AC-022.1: Cat\'alogo accesible en modo avi\'on. \newline
AC-022.2: Im\'agenes visibles sin conexi\'on. \\
\hline
\end{longtable}

\begin{longtable}{|p{3cm}|p{11.5cm}|}
\hline
\multicolumn{2}{|c|}{\textbf{RF-023}} \\
\hline
Nombre & Edici\'on de fichas del cat\'alogo (Bot\'anico/Administrador -- Web) \\
\hline
Obligatoriedad & \textbf{deber\'a} \\
\hline
Descripci\'on & La plataforma web deber\'a permitir al Bot\'anico y al Administrador agregar, editar y marcar como validada cada ficha del cat\'alogo, incluyendo los 11 atributos Ochoa y los atributos visuales adicionales. Toda edici\'on registrada con usuario y fecha. \\
\hline
Fuente & CRISP-DM Fase 3. \\
\hline
Prioridad & Alta \\
\hline
Criterios de Aceptaci\'on &
AC-023.1: Bot\'anico edita todos los campos y guarda. \newline
AC-023.2: Historial de ediciones con qui\'en y cu\'ando. \\
\hline
\end{longtable}

%-------------------------------------------
\subsection{M\'odulo 5: Registro de Observaciones en Campo}
%-------------------------------------------

\textit{Plataforma: App m\'ovil exclusivamente.}

\begin{longtable}{|p{3cm}|p{11.5cm}|}
\hline
\multicolumn{2}{|c|}{\textbf{RF-024} \criticomark} \\
\hline
Nombre & Crear registro de observaci\'on con informaci\'on biol\'ogica completa \\
\hline
Obligatoriedad & \textbf{deber\'a} \\
\hline
Descripci\'on & La app deber\'a crear un registro de observaci\'on con:\newline
\textbf{Autom\'aticos:} especie identificada (o ``no identificada''), fecha/hora, coordenadas GPS, fotograf\'ia(s).\newline
\textbf{Requeridos:} estado del individuo (floraci\'on/crecimiento/marchito), puesto de control, notas libres.\newline
\textbf{Informaci\'on biol\'ogica adicional} (opcional pero solicitada): color del labelo observado, presencia de flores abiertas, condiciones clim\'aticas (sol/nublado/niebla/lluvia), microh\'abitat (\'arbol/roca/suelo). \\
\hline
Rationale & La informaci\'on biol\'ogica de campo (estado fenol\'ogico, microh\'abitat, condiciones) aumenta el valor cient\'ifico del registro y facilita la validaci\'on por el Validador. \\
\hline
Fuente & Bot\'anico especialista. ANEXO 2. \\
\hline
Prioridad & Alta \\
\hline
Verificaci\'on & Prueba funcional en campo, incluyendo escenarios con se\~nal GPS d\'ebil. \\
\hline
Stakeholder & Guardaparque \\
\hline
Criterios de Aceptaci\'on &
AC-024.1: Registro creado en $\leq$30 segundos tras la identificaci\'on. \newline
AC-024.2: GPS: intento $\leq$30s; si no alcanza $\leq$10m, aplica degradaci\'on elegante: guarda mejor ubicaci\'on disponible con margen real en campo \texttt{precision\_gps\_m}, mensaje ``Se\~nal GPS d\'ebil -- guardando ubicaci\'on aproximada ($\pm$Xm)''. \newline
AC-024.3: Registro guardado localmente sin conexi\'on. \newline
AC-024.4: Campo \texttt{datos\_campo\_completos} (bool) indica si se proporcionaron los campos biol\'ogicos opcionales. \\
\hline
\end{longtable}

\begin{longtable}{|p{3cm}|p{11.5cm}|}
\hline
\multicolumn{2}{|c|}{\textbf{RF-025} \criticomark} \\
\hline
Nombre & Almacenamiento local con pol\'itica de purga \\
\hline
Obligatoriedad & \textbf{deber\'a} \\
\hline
Descripci\'on & La app deber\'a almacenar observaciones en SQLite/Room local. Una vez sincronizada y validada, elimina las fotos originales de alta resoluci\'on y conserva solo metadatos y thumbnail ($\leq$200$\times$200 px). \\
\hline
Fuente & CR-03. \\
\hline
Prioridad & Alta \\
\hline
Criterios de Aceptaci\'on &
AC-025.1: Cierre forzado durante guardado: registro conservado al reiniciar. \newline
AC-025.2: BD local $\geq$10,000 registros sin degradaci\'on. \newline
AC-025.3: Fotos purgadas post-sincronizaci\'on+validaci\'on; thumbnail conservado. \newline
AC-025.4: Thumbnail visible en historial de observaciones purgadas. \\
\hline
\end{longtable}

\begin{longtable}{|p{3cm}|p{11.5cm}|}
\hline
\multicolumn{2}{|c|}{\textbf{RF-026} \criticomark} \\
\hline
Nombre & Sincronizaci\'on autom\'atica diferida \\
\hline
Obligatoriedad & \textbf{deber\'a} \\
\hline
Descripci\'on & La app deber\'a sincronizar observaciones pendientes al recuperar conexi\'on, sin intervenci\'on del usuario. Reintento con retroceso exponencial (m\'ax. 5). Pol\'itica de merge: campos de validaci\'on desde servidor (fuente de verdad); campo \texttt{notas} concatenado. Notificaci\'on: ``Tu observaci\'on fue actualizada. Tus notas han sido conservadas.'' \\
\hline
Fuente & CRISP-DM Fase 6. \\
\hline
Prioridad & Alta \\
\hline
Criterios de Aceptaci\'on &
AC-026.1: Sincronizaci\'on en $\leq$2 minutos al recuperar conexi\'on. \newline
AC-026.2: Sin duplicados en servidor. \newline
AC-026.3: Conflicto: campos validaci\'on desde servidor; \texttt{notas} concatenado. \newline
AC-026.4: Notificaci\'on enviada al Guardaparque. \\
\hline
\end{longtable}

% --- NUEVO RF-060 (Cola de sincronización asíncrona) ---
\begin{longtable}{|p{3cm}|p{11.5cm}|}
\hline
\multicolumn{2}{|c|}{\textbf{RF-060} \criticomark} \\
\hline
Nombre & Cola de sincronización asíncrona de registros \\
\hline
Obligatoriedad & \textbf{deberá} \\
\hline
Descripción & El sistema debe implementar una cola de sincronización local. Los registros creados en campo se marcarán como \texttt{‘Pendientes’}. Una vez que el dispositivo detecte una conexión estable (WiFi/4G), el sistema deberá subir los datos y las imágenes al API Backend y, opcionalmente, enviar la imagen a Hugging Face/Servidor Central solo para procesos de re-entrenamiento o validación por expertos. \\
\hline
Rationale & Asegura la entrega ordenada de observaciones y permite enriquecer el dataset futuro sin depender de sincronización manual. \\
\hline
Fuente & Arquitectura Offline-First. \\
\hline
Prioridad & Alta \\
\hline
Verificación & Interrumpir la red durante la sincronización; verificar que los registros quedan en cola y se reintentan automáticamente. \\
\hline
Stakeholder & Guardaparque, Administrador \\
\hline
Criterios de Aceptación &
AC-060.1: Cada registro local tiene un estado de sincronización (pendiente/enviado). \newline
AC-060.2: Al conectar, los registros pendientes se envían en orden cronológico ascendente. \newline
AC-060.3: Si la sincronización se interrumpe, el sistema reintenta al recuperar la conexión. \newline
AC-060.4: Solo cuando el registro se confirma en el servidor, se marca como enviado. \newline
AC-060.5: El envío opcional a servidor de reentrenamiento no bloquea la sincronización principal. \\
\hline
\end{longtable}

\begin{longtable}{|p{3cm}|p{11.5cm}|}
\hline
\multicolumn{2}{|c|}{\textbf{RF-027}} \\
\hline
Nombre & Edici\'on de observaciones propias \\
\hline
Obligatoriedad & \textbf{deber\'ia} \\
\hline
Descripci\'on & La app deber\'ia permitir al Guardaparque editar sus propios registros no validados: especie, notas, estado del individuo y datos biol\'ogicos de campo. \\
\hline
Fuente & \inferido{} \\
\hline
Prioridad & Media \\
\hline
Criterios de Aceptaci\'on &
AC-027.1: Guardaparque edita observaci\'on no validada. \newline
AC-027.2: Observaci\'on validada no puede ser editada por el Guardaparque. \\
\hline
\end{longtable}

\begin{longtable}{|p{3cm}|p{11.5cm}|}
\hline
\multicolumn{2}{|c|}{\textbf{RF-028}} \\
\hline
Nombre & Historial de observaciones del Guardaparque \\
\hline
Obligatoriedad & \textbf{deber\'a} \\
\hline
Descripci\'on & La app deber\'a mostrar al Guardaparque su historial con filtros por fecha, especie y estado de validaci\'on. Observaciones purgadas muestran thumbnail. \\
\hline
Fuente & \inferido{} \\
\hline
Prioridad & Media \\
\hline
Criterios de Aceptaci\'on &
AC-028.1: Historial ordenado por fecha descendente. \newline
AC-028.2: Filtro por estado de validaci\'on funcional (pendiente/validada/corregida). \\
\hline
\end{longtable}

\begin{longtable}{|p{3cm}|p{11.5cm}|}
\hline
\multicolumn{2}{|c|}{\textbf{RF-029}} \\
\hline
Nombre & Exportaci\'on de registros a CSV \\
\hline
Obligatoriedad & \textbf{deber\'ia} \\
\hline
Descripci\'on & El sistema deber\'ia permitir exportar observaciones filtradas en CSV (app: Guardaparque; web: Validador/Administrador), incluyendo \texttt{precision\_gps\_m} y datos biol\'ogicos de campo. \\
\hline
Fuente & \inferido{} \\
\hline
Prioridad & Media \\
\hline
Criterios de Aceptaci\'on &
AC-029.1: CSV incluye todos los campos del registro. \newline
AC-029.2: Archivo guardado en almacenamiento accesible al usuario. \\
\hline
\end{longtable}

%-------------------------------------------
\subsection{M\'odulo 6: Validaci\'on por el Validador}
%-------------------------------------------

\textit{Plataforma: Plataforma web exclusivamente (Actor 2 -- Validador).}

\begin{longtable}{|p{3cm}|p{11.5cm}|}
\hline
\multicolumn{2}{|c|}{\textbf{RF-030} \criticomark} \\
\hline
Nombre & Dashboard web de observaciones pendientes de validaci\'on \\
\hline
Obligatoriedad & \textbf{deber\'a} \\
\hline
Descripci\'on & La plataforma web deber\'a mostrar al Validador un dashboard con todas las observaciones pendientes de validaci\'on, con: miniatura de foto, especie identificada por la CNN, nombre del Guardaparque, puesto de control, fecha/hora, coordenadas GPS y margen de error, estado fenol\'ogico, datos biol\'ogicos de campo, respuestas al flujo de preguntas guiadas (si aplica). Ordenadas por fecha; filtros por puesto, especie, estado y tipo (normal / flujo asistido / no identificada). \\
\hline
Rationale & El Validador revisa el trabajo de 40 guardaparques desde su estaci\'on de trabajo web. Necesita toda la informaci\'on biol\'ogica en una sola vista para tomar decisiones informadas. \\
\hline
Fuente & ANEXO 2 (sobrecarga de validaci\'on). Ref. 10. \\
\hline
Prioridad & Alta \\
\hline
Verificaci\'on & Prueba funcional web con 20 observaciones de distintos tipos. \\
\hline
Stakeholder & Validador \\
\hline
Criterios de Aceptaci\'on &
AC-030.1: Dashboard muestra todos los campos por observaci\'on. \newline
AC-030.2: Filtros por puesto, especie, estado y tipo funcionales. \newline
AC-030.3: Lista actualizada al sincronizarse nuevas observaciones. \\
\hline
\end{longtable}

\begin{longtable}{|p{3cm}|p{11.5cm}|}
\hline
\multicolumn{2}{|c|}{\textbf{RF-031} \criticomark} \\
\hline
Nombre & Aprobar, corregir o rechazar identificaciones (Web) \\
\hline
Obligatoriedad & \textbf{deber\'a} \\
\hline
Descripci\'on & La plataforma web deber\'a permitir al Validador: aprobar, corregir (seleccionando especie del cat\'alogo de 16) o rechazar (con justificaci\'on $\geq$10 caracteres) cada observaci\'on. Toda acci\'on registrada con timestamp y usuario. \\
\hline
Fuente & ANEXO 2. CRISP-DM Fase 5. \\
\hline
Prioridad & Alta \\
\hline
Criterios de Aceptaci\'on &
AC-031.1: Aprobaci\'on: marca ``Validado'' con nombre del Validador y fecha. \newline
AC-031.2: Correcci\'on: registra especie original del Guardaparque y especie corregida. \newline
AC-031.3: Rechazo: justificaci\'on obligatoria $\geq$10 caracteres. \\
\hline
\end{longtable}

\begin{longtable}{|p{3cm}|p{11.5cm}|}
\hline
\multicolumn{2}{|c|}{\textbf{RF-032}} \\
\hline
Nombre & Comentarios t\'ecnicos por observaci\'on (Web) \\
\hline
Obligatoriedad & \textbf{deber\'ia} \\
\hline
Descripci\'on & La plataforma web deber\'ia permitir al Validador a\~nadir comentarios t\'ecnicos sobre caracter\'isticas morfol\'ogicas (labelo, florecimiento, hojas) para retroalimentar al Guardaparque. \\
\hline
Fuente & ANEXO 2. \\
\hline
Prioridad & Media \\
\hline
Criterios de Aceptaci\'on &
AC-032.1: Guardaparque ve el comentario en su historial de la app. \\
\hline
\end{longtable}

\begin{longtable}{|p{3cm}|p{11.5cm}|}
\hline
\multicolumn{2}{|c|}{\textbf{RF-033}} \\
\hline
Nombre & Notificaci\'on de resultado de validaci\'on (App) \\
\hline
Obligatoriedad & \textbf{deber\'ia} \\
\hline
Descripci\'on & La app deber\'ia notificar al Guardaparque cuando una observaci\'on haya sido validada, corregida o rechazada desde la plataforma web. \\
\hline
Fuente & \inferido{} \\
\hline
Prioridad & Media \\
\hline
Criterios de Aceptaci\'on &
AC-033.1: Al sincronizarse, Guardaparque recibe notificaci\'on con resultado. \\
\hline
\end{longtable}

\begin{longtable}{|p{3cm}|p{11.5cm}|}
\hline
\multicolumn{2}{|c|}{\textbf{RF-034}} \\
\hline
Nombre & Estad\'isticas de precisi\'on por Guardaparque (Web) \\
\hline
Obligatoriedad & \textbf{deber\'ia} \\
\hline
Descripci\'on & La plataforma web deber\'ia mostrar al Validador estad\'isticas de precisi\'on por Guardaparque: tasa de acierto, especies con m\'as errores, tasa de activaci\'on del flujo asistido y tendencia temporal. \\
\hline
Fuente & CRISP-DM Fase 5. \\
\hline
Prioridad & Media \\
\hline
Criterios de Aceptaci\'on &
AC-034.1: Tasa de acierto por Guardaparque en per\'iodo seleccionado. \newline
AC-034.2: Las 3 especies con mayor tasa de error por Guardaparque. \\
\hline
\end{longtable}

%-------------------------------------------
\subsection{M\'odulo 7: Mapa de Observaciones}
%-------------------------------------------

\textit{Plataforma: App m\'ovil (visualizaci\'on b\'asica offline) + Plataforma web (mapa interactivo con zonas de inter\'es).}

\begin{longtable}{|p{3cm}|p{11.5cm}|}
\hline
\multicolumn{2}{|c|}{\textbf{RF-035} \criticomark} \\
\hline
Nombre & Mapa de observaciones georreferenciadas (App m\'ovil) \\
\hline
Obligatoriedad & \textbf{deber\'a} \\
\hline
Descripci\'on & La app m\'ovil deber\'a mostrar un mapa del SHMP con:\newline
\textbf{(a)} Puntos georreferenciados de las observaciones del Guardaparque, con color diferenciado por estado de validaci\'on (pendiente: gris, validado: verde, corregido: azul, rechazado: rojo).\newline
\textbf{(b)} Puntos con \texttt{precision\_gps\_m} $>$10m muestran radio de incertidumbre.\newline
\textbf{(c)} Capa de \textbf{zonas de inter\'es del SHMP}: sectores principales (Aguas Calientes, Camino Inca, zona arqueol\'ogica, bosque montano) como pol\'igonos informativos que el Guardaparque puede consultar para orientarse. \\
\hline
Rationale & El mapa permite al Guardaparque visualizar su trabajo de campo en contexto geogr\'afico, identificar zonas de alta densidad de observaciones propias y orientarse en el terreno. \\
\hline
Fuente & ANEXO 2. Bot\'anico especialista. \\
\hline
Prioridad & Alta \\
\hline
Verificaci\'on & Prueba funcional con 20 observaciones georreferenciadas. \\
\hline
Stakeholder & Guardaparque \\
\hline
Criterios de Aceptaci\'on &
AC-035.1: Puntos visibles con diferenciaci\'on de color por estado de validaci\'on. \newline
AC-035.2: Tap en punto: resumen de la observaci\'on (especie, fecha, estado). \newline
AC-035.3: Puntos de baja precisi\'on GPS muestran radio de incertidumbre. \newline
AC-035.4: Capa de zonas de inter\'es visible como capa activable/desactivable. \\
\hline
\end{longtable}

\begin{longtable}{|p{3cm}|p{11.5cm}|}
\hline
\multicolumn{2}{|c|}{\textbf{RF-052} \criticomark} \\
\hline
Nombre & Mapa interactivo de observaciones con zonas de inter\'es (Plataforma web) \\
\hline
Obligatoriedad & \textbf{deber\'a} \\
\hline
Descripci\'on & La plataforma web deber\'a mostrar al Validador y al Administrador un mapa interactivo del SHMP con:\newline
\textbf{(a)} Todos los puntos de observaci\'on georreferenciados con diferenciaci\'on de color por estado de validaci\'on y especie.\newline
\textbf{(b)} Capas de \textbf{zonas de inter\'es}: sectores del Santuario, los 7 puestos de control del SERNANP (con nombre y guardaparques asignados), rutas de patrullaje principales.\newline
\textbf{(c)} Filtros por especie, rango de fechas, puesto de control y estado de validaci\'on.\newline
\textbf{(d)} Al seleccionar un punto: vista detallada de la observaci\'on con fotograf\'ia, datos de campo y estado de validaci\'on. \\
\hline
Rationale & El Validador necesita visualizar la distribuci\'on geogr\'afica de las observaciones en contexto del Santuario, los puestos de control y las rutas, para tomar decisiones sobre la cobertura del monitoreo. \\
\hline
Fuente & ANEXO 2. Ref. 10. \\
\hline
Prioridad & Alta \\
\hline
Verificaci\'on & Prueba funcional web con $\geq$50 observaciones de distintas zonas. \\
\hline
Stakeholder & Validador, Administrador \\
\hline
Criterios de Aceptaci\'on &
AC-052.1: Mapa renderiza en $\leq$5 segundos con $\leq$500 observaciones. \newline
AC-052.2: Puntos con diferenciaci\'on de color por estado y especie. \newline
AC-052.3: Capa de puestos de control visible y activable independientemente. \newline
AC-052.4: Filtros de especie y per\'iodo actualizan el mapa din\'amicamente. \newline
AC-052.5: Clic en punto: vista detallada de la observaci\'on. \\
\hline
\end{longtable}

\begin{longtable}{|p{3cm}|p{11.5cm}|}
\hline
\multicolumn{2}{|c|}{\textbf{RF-036}} \\
\hline
Nombre & Filtrado de observaciones en mapa (App y Web) \\
\hline
Obligatoriedad & \textbf{deber\'ia} \\
\hline
Descripci\'on & El mapa (app y web) deber\'ia permitir filtrar por especie, rango de fechas, puesto de control y estado de validaci\'on. \\
\hline
Fuente & \inferido{} \\
\hline
Prioridad & Media \\
\hline
Criterios de Aceptaci\'on &
AC-036.1: Filtro por especie muestra solo observaciones de esa especie. \\
\hline
\end{longtable}

\begin{longtable}{|p{3cm}|p{11.5cm}|}
\hline
\multicolumn{2}{|c|}{\textbf{RF-037} \criticomark} \\
\hline
Nombre & Georreferenciaci\'on autom\'atica con degradaci\'on elegante \\
\hline
Obligatoriedad & \textbf{deber\'a} \\
\hline
Descripci\'on & La app deber\'a obtener coordenadas GPS autom\'aticamente al crear una observaci\'on. Si en 30s no alcanza $\leq$10m de precisi\'on, aplica degradaci\'on elegante (AC-024.2). GPS activado solo bajo demanda (RNF-020). \\
\hline
Fuente & ANEXO 2. \\
\hline
Prioridad & Alta \\
\hline
Criterios de Aceptaci\'on &
AC-037.1: GPS intentado en $\leq$30s. \newline
AC-037.2: GPS no disponible: notifica y permite continuar sin coordenadas. \newline
AC-037.3: Precisi\'on $>$10m antes de 30s: aplica degradaci\'on elegante de AC-024.2. \\
\hline
\end{longtable}

\begin{longtable}{|p{3cm}|p{11.5cm}|}
\hline
\multicolumn{2}{|c|}{\textbf{RF-038}} \\
\hline
Nombre & Mapa b\'asico offline con l\'imite de zoom (App) \\
\hline
Obligatoriedad & \textbf{deber\'ia} \\
\hline
Descripci\'on & La app deber\'ia almacenar en cach\'e los tiles del mapa del SHMP hasta \textbf{Zoom Level 15/16 m\'aximo} ($\leq$50 MB) para visualizaci\'on offline. Zoom superior muestra mensaje informativo. \\
\hline
Fuente & CR-04. \\
\hline
Prioridad & Media \\
\hline
Criterios de Aceptaci\'on &
AC-038.1: Mapa visible offline con observaciones locales. \newline
AC-038.2: Cach\'e no supera ZL 16. \newline
AC-038.3: Peso cach\'e $\leq$50 MB. \newline
AC-038.4: Zoom $>$16 offline: mensaje informativo. \\
\hline
\end{longtable}

%-------------------------------------------
\subsection{M\'odulo 8: Reportes e Informes}
%-------------------------------------------

\textit{Plataforma: Plataforma web (generaci\'on y exportaci\'on) + App m\'ovil (encuesta de usabilidad).}

\begin{longtable}{|p{3cm}|p{11.5cm}|}
\hline
\multicolumn{2}{|c|}{\textbf{RF-039}} \\
\hline
Nombre & Generaci\'on de informe de monitoreo (Web) \\
\hline
Obligatoriedad & \textbf{deber\'ia} \\
\hline
Descripci\'on & La plataforma web deber\'ia generar informes por puesto de control y per\'iodo: especies observadas, cantidad de registros, tasa de identificaci\'on correcta, tasa de activaci\'on del flujo asistido, reportes de especie no identificada y mapa de puntos. \\
\hline
Fuente & ANEXO 2. \\
\hline
Prioridad & Alta \\
\hline
Criterios de Aceptaci\'on &
AC-039.1: Informe generado en $\leq$10 segundos con $\leq$200 observaciones. \newline
AC-039.2: Incluye todos los campos definidos. \\
\hline
\end{longtable}

\begin{longtable}{|p{3cm}|p{11.5cm}|}
\hline
\multicolumn{2}{|c|}{\textbf{RF-040}} \\
\hline
Nombre & Exportaci\'on de informes en PDF (Web) \\
\hline
Obligatoriedad & \textbf{deber\'ia} \\
\hline
Descripci\'on & La plataforma web deber\'ia exportar informes en PDF con datos georreferenciados para reportes institucionales del SERNANP. \\
\hline
Fuente & ANEXO 2. \\
\hline
Prioridad & Alta \\
\hline
Criterios de Aceptaci\'on &
AC-040.1: PDF v\'alido, legible en cualquier visor. \newline
AC-040.2: PDF incluye mapa del SHMP con puntos del per\'iodo. \\
\hline
\end{longtable}

\begin{longtable}{|p{3cm}|p{11.5cm}|}
\hline
\multicolumn{2}{|c|}{\textbf{RF-041}} \\
\hline
Nombre & M\'etricas de uso y precisi\'on del sistema (Web) \\
\hline
Obligatoriedad & \textbf{deber\'ia} \\
\hline
Descripci\'on & La plataforma web deber\'ia mostrar al Validador y Administrador: tasa de acierto global del modelo, tasa de activaci\'on del flujo asistido por especie, tiempo promedio de identificaci\'on, n\'umero de reportes de especie desconocida y tendencia mensual. \\
\hline
Fuente & CRISP-DM Fase 5. \\
\hline
Prioridad & Media \\
\hline
Criterios de Aceptaci\'on &
AC-041.1: Panel muestra tasa de acierto global y tiempo promedio. \newline
AC-041.2: Panel muestra tasa de activaci\'on del flujo asistido por especie. \\
\hline
\end{longtable}

\begin{longtable}{|p{3cm}|p{11.5cm}|}
\hline
\multicolumn{2}{|c|}{\textbf{RF-042}} \\
\hline
Nombre & Encuesta de usabilidad (App) \\
\hline
Obligatoriedad & \textbf{deber\'ia} \\
\hline
Descripci\'on & La app deber\'ia mostrar encuesta corta (5 preguntas Likert 1--5) una vez por d\'ia al primer cierre de sesi\'on, si se registraron $\geq$3 observaciones esa jornada. Resultados agrupados en panel web del Administrador. \\
\hline
Fuente & ANEXO 2 (Objetivo espec\'ifico 5). \\
\hline
Prioridad & Media \\
\hline
Criterios de Aceptaci\'on &
AC-042.1: Encuesta aparece en jornadas con $\geq$3 observaciones. \newline
AC-042.2: Resultados visibles en panel del Administrador. \\
\hline
\end{longtable}

%-------------------------------------------
\subsection{M\'odulo 9: Administraci\'on del Sistema}
%-------------------------------------------

\textit{Plataforma: Plataforma web exclusivamente (Actor 4 -- Administrador).}

\begin{longtable}{|p{3cm}|p{11.5cm}|}
\hline
\multicolumn{2}{|c|}{\textbf{RF-043} \criticomark} \\
\hline
Nombre & Gesti\'on de usuarios CRUD (Administrador -- Web) \\
\hline
Obligatoriedad & \textbf{deber\'a} \\
\hline
Descripci\'on & La plataforma web deber\'a permitir al Administrador (\'unico actor autorizado) crear, leer, actualizar y desactivar cuentas de usuario para los 4 roles, asignar y cambiar roles, y reasignar puestos de control para los 40 guardaparques en 7 puestos del SHMP. Los datos del formulario de creaci\'on se env\'ian en JSON al API Backend (HTTP 201 -- Ref. 10, RF-01). \\
\hline
Fuente & ANEXO 2 (40 guardaparques, 7 puestos). Ref. 10. \\
\hline
Prioridad & Alta \\
\hline
Verificaci\'on & Prueba de cada operaci\'on CRUD desde la web + verificaci\'on en BD PostgreSQL. \\
\hline
Stakeholder & Administrador \\
\hline
Criterios de Aceptaci\'on &
AC-043.1: Administrador crea, edita y desactiva usuarios desde la web. \newline
AC-043.2: Usuario desactivado no puede iniciar sesi\'on en ninguna plataforma. \newline
AC-043.3: Cambio de rol aplicado en pr\'oximo inicio de sesi\'on del usuario. \newline
AC-043.4: Contrase\~nas almacenadas con hash Bcrypt en PostgreSQL (no texto plano -- Ref. 10, RNF-03). \\
\hline
\end{longtable}

\begin{longtable}{|p{3cm}|p{11.5cm}|}
\hline
\multicolumn{2}{|c|}{\textbf{RF-044} \criticomark} \\
\hline
Nombre & Actualizaci\'on OTA del modelo CNN con reanudaci\'on de descarga \\
\hline
Obligatoriedad & \textbf{deber\'a} \\
\hline
Descripci\'on & La plataforma web deber\'a permitir al Administrador subir nueva versi\'on del modelo .tflite. La app deber\'a descargarlo y aplicarlo sin reinstalaci\'on. Modelo anterior conservado como respaldo. La descarga OTA soporta reanudaci\'on ante cortes de red (CR-05). \\
\hline
Fuente & CR-05. CRISP-DM Fase 6. \\
\hline
Prioridad & Alta \\
\hline
Criterios de Aceptaci\'on &
AC-044.1: Administrador sube .tflite desde panel web. \newline
AC-044.2: Dispositivos descargan nuevo modelo al conectarse. \newline
AC-044.3: Versi\'on anterior almacenada como respaldo local. \newline
AC-044.4: Descarga interrumpida reanuda desde el punto de interrupci\'on. \newline
AC-044.5: Descarga verificada con checksum SHA-256 activa autom\'aticamente el nuevo modelo. \\
\hline
\end{longtable}

\begin{longtable}{|p{3cm}|p{11.5cm}|}
\hline
\multicolumn{2}{|c|}{\textbf{RF-045}} \\
\hline
Nombre & Gesti\'on del cat\'alogo de 16 especies (Web) \\
\hline
Obligatoriedad & \textbf{deber\'a} \\
\hline
Descripci\'on & La plataforma web deber\'a permitir al Administrador agregar nuevas especies (con 11 atributos Ochoa + atributos visuales), editar fichas y actualizar el cat\'alogo en todos los dispositivos mediante sincronizaci\'on. \\
\hline
Fuente & CRISP-DM Fase 6. \\
\hline
Prioridad & Alta \\
\hline
Criterios de Aceptaci\'on &
AC-045.1: Nueva especie a\~nadida desde web aparece en cat\'alogo de todos los dispositivos tras sincronizaci\'on. \\
\hline
\end{longtable}

\begin{longtable}{|p{3cm}|p{11.5cm}|}
\hline
\multicolumn{2}{|c|}{\textbf{RF-046}} \\
\hline
Nombre & Panel de estad\'isticas globales del sistema (Administrador -- Web) \\
\hline
Obligatoriedad & \textbf{deber\'ia} \\
\hline
Descripci\'on & La plataforma web deber\'ia mostrar al Administrador: usuarios activos por rol, observaciones totales y pendientes de validaci\'on, tasa de sincronizaci\'on, versi\'on del modelo en uso, porcentaje de almacenamiento por dispositivo, reportes de especie desconocida pendientes y errores del sistema o del API Backend. \\
\hline
Fuente & \inferido{}. Ref. 10 (monitoreo del API). \\
\hline
Prioridad & Media \\
\hline
Criterios de Aceptaci\'on &
AC-046.1: Panel muestra en tiempo real usuarios activos y observaciones pendientes de sincronizaci\'on. \newline
AC-046.2: Panel muestra reportes de especie desconocida pendientes de revisi\'on. \newline
AC-046.3: Versi\'on del modelo activo visible. \\
\hline
\end{longtable}

\begin{longtable}{|p{3cm}|p{11.5cm}|}
\hline
\multicolumn{2}{|c|}{\textbf{RF-047} \criticomark} \\
\hline
Nombre & Consentimiento de privacidad en primer uso \\
\hline
Obligatoriedad & \textbf{deber\'a} \\
\hline
Descripci\'on & La app deber\'a mostrar pantalla de consentimiento informado (Ley N.° 29733) en el primer inicio. La plataforma web deber\'a mostrar aviso de privacidad en el primer acceso. Sin aceptaci\'on: acceso bloqueado en ambas plataformas. \\
\hline
Fuente & ANEXO 2 (Componente \'Etico). \\
\hline
Prioridad & Alta \\
\hline
Criterios de Aceptaci\'on &
AC-047.1: App primer inicio: consentimiento antes de cualquier funcionalidad. \newline
AC-047.2: Web primer acceso: aviso de privacidad. \newline
AC-047.3: Sin aceptaci\'on: acceso bloqueado. \\
\hline
\end{longtable}

\begin{longtable}{|p{3cm}|p{11.5cm}|}
\hline
\multicolumn{2}{|c|}{\textbf{RF-048}} \\
\hline
Nombre & Manual de usuario integrado \\
\hline
Obligatoriedad & \textbf{deber\'ia} \\
\hline
Descripci\'on & La app deber\'ia incluir manual integrado offline con gu\'ias del flujo principal y del flujo asistido. La plataforma web deber\'ia incluir secci\'on de ayuda en l\'inea. \\
\hline
Fuente & ANEXO 2 (Resultado esperado 7). \\
\hline
Prioridad & Media \\
\hline
Criterios de Aceptaci\'on &
AC-048.1: Manual app accesible offline desde men\'u de ayuda. \newline
AC-048.2: Ayuda web accesible desde navegador. \\
\hline
\end{longtable}

%-------------------------------------------
\subsection{M\'odulo 10: Revisi\'on de Reportes de Especie Desconocida (Bot\'anico -- Web)}
%-------------------------------------------

\textit{Plataforma: Plataforma web exclusivamente (Actor 3 -- Bot\'anico).}

\begin{longtable}{|p{3cm}|p{11.5cm}|}
\hline
\multicolumn{2}{|c|}{\textbf{RF-055} \criticomark} \\
\hline
Nombre & Cola web de reportes de especie no identificada (Bot\'anico) \\
\hline
Obligatoriedad & \textbf{deber\'a} \\
\hline
Descripci\'on & La plataforma web deber\'a proveer al Bot\'anico una cola de revisi\'on de reportes de especie no identificada (RF-051). Para cada reporte, el Bot\'anico ver\'a: todas las fotograf\'ias de m\'ultiples \'angulos, respuestas del Guardaparque a las preguntas guiadas, coordenadas GPS, fecha/hora, estado fenol\'ogico, datos de campo y vector de probabilidades del \'ultimo intento CNN. El Bot\'anico podr\'a:\newline
(a) Asignar el reporte a una de las 16 especies del cat\'alogo.\newline
(b) Descartar el reporte (imagen insuficiente para determinar la especie).\newline
(c) Marcar como ``Posible nueva especie -- requiere an\'alisis ex-situ''. \\
\hline
Rationale & El Bot\'anico es el \'unico actor capacitado para determinar si un reporte corresponde a una especie ya catalogada, una variante morfol\'ogica o una especie potencialmente nueva para el SHMP. \\
\hline
Fuente & Bot\'anico especialista. RF-051. \\
\hline
Prioridad & Alta \\
\hline
Verificaci\'on & Prueba con cola de 5 reportes de distintos tipos. \\
\hline
Stakeholder & Bot\'anico \\
\hline
Criterios de Aceptaci\'on &
AC-055.1: Cola ordenada por prioridad (alertas ``Posible Nueva Especie'' primero). \newline
AC-055.2: Vista detalle: todas las im\'agenes, respuestas a preguntas y vector CNN. \newline
AC-055.3: Bot\'anico puede asignar especie, descartar o marcar como nueva; acci\'on registrada con timestamp. \\
\hline
\end{longtable}

\begin{longtable}{|p{3cm}|p{11.5cm}|}
\hline
\multicolumn{2}{|c|}{\textbf{RF-056}} \\
\hline
Nombre & Tendencias temporales por especie (Bot\'anico/Validador -- Web) \\
\hline
Obligatoriedad & \textbf{deber\'ia} \\
\hline
Descripci\'on & La plataforma web deber\'ia mostrar gr\'aficos de tendencia temporal de observaciones por especie: distribuci\'on mensual de avistamientos y \'epoca de floraci\'on observada vs. esperada seg\'un ficha Ochoa. \\
\hline
Fuente & Bot\'anico especialista. RNF-019. \\
\hline
Prioridad & Media \\
\hline
Criterios de Aceptaci\'on &
AC-056.1: Gr\'afico mensual por especie renderiza correctamente. \newline
AC-056.2: \'Epoca observada vs. esperada visible en el mismo gr\'afico. \\
\hline
\end{longtable}

%-------------------------------------------
\subsection{RF relacionados con el API Backend}
%-------------------------------------------

\textit{Los siguientes RF derivan directamente del documento de arquitectura del equipo web (Ref. 10) y formalizan los requerimientos de la capa de integraci\'on.}

\begin{longtable}{|p{3cm}|p{11.5cm}|}
\hline
\multicolumn{2}{|c|}{\textbf{RF-057} \criticomark} \\
\hline
Nombre & Compatibilidad estricta del API con la App M\'ovil \\
\hline
Obligatoriedad & \textbf{deber\'a} \\
\hline
Descripci\'on & El API Backend (Node.js + Express) deber\'a exponer \textit{endpoints} con estructuras de datos JSON id\'enticas a las que la app m\'ovil espera recibir, de modo que el cliente m\'ovil funcione sin detectar cambios. Cualquier nueva funcionalidad del API se agrega como \textit{endpoint} adicional sin modificar los existentes. \\
\hline
Rationale & El c\'odigo de la app m\'ovil es el cliente base del sistema. El dise\~no de las respuestas del API est\'a dictado por lo que la app necesita (Ref. 10, secci\'on 6). \\
\hline
Fuente & Ref. 10 (RF-02, secci\'on 3.2). \\
\hline
Prioridad & Alta \\
\hline
Verificaci\'on & Pruebas de integraci\'on: peticiones id\'enticas a las de la app m\'ovil al API retornan respuestas 100\% compatibles. \\
\hline
Stakeholder & Equipo de desarrollo web y m\'ovil \\
\hline
Criterios de Aceptaci\'on &
AC-057.1: Petici\'on de la app m\'ovil al nuevo API: respuesta id\'entica a la esperada por el cliente m\'ovil. \newline
AC-057.2: No se modifican \textit{endpoints} existentes al a\~nadir nueva funcionalidad. \\
\hline
\end{longtable}

\begin{longtable}{|p{3cm}|p{11.5cm}|}
\hline
\multicolumn{2}{|c|}{\textbf{RF-058} \criticomark} \\
\hline
Nombre & Comunicaci\'on React $\leftrightarrow$ API en formato JSON \\
\hline
Obligatoriedad & \textbf{deber\'a} \\
\hline
Descripci\'on & La plataforma web (React.js) deber\'a enviar todos los datos de sus formularios al API Backend en formato JSON mediante Axios/Fetch API. El API deber\'a retornar HTTP 201 (Creado) ante peticiones POST exitosas de creaci\'on de recursos. \\
\hline
Fuente & Ref. 10 (RF-01). \\
\hline
Prioridad & Alta \\
\hline
Verificaci\'on & Inspecci\'on de red (DevTools): petici\'on POST desde React retorna HTTP 201. \\
\hline
Criterios de Aceptaci\'on &
AC-058.1: POST de formulario React: API retorna HTTP 201 con el recurso creado. \newline
AC-058.2: Validaci\'on en tiempo real en el navegador antes de enviar (campos obligatorios, formato email). \\
\hline
\end{longtable}

%=============================================================
\section{Requerimientos no funcionales}
%=============================================================

Clasificados seg\'un ISO/IEC 25010:2011.

\subsection{Desempe\~no y eficiencia}

\begin{longtable}{|p{3cm}|p{11.5cm}|}
\hline
\multicolumn{2}{|c|}{\textbf{RNF-001} \criticomark} \\
\hline
Categor\'ia & Eficiencia de desempe\~no -- App m\'ovil \\
\hline
Obligatoriedad & \textbf{deber\'a} \\
\hline
Descripci\'on & Tiempo total de identificaci\'on (captura confirmada $\rightarrow$ resultados) $\leq$5 segundos en dispositivo con 2 GB RAM y procesador octa-core 1.8 GHz. \\
\hline
Fuente & ANEXO 2 (Hip\'otesis 2). \\
\hline
Prioridad & Alta \\
\hline
Verificaci\'on & Prueba de rendimiento con 10 iteraciones; registrar promedio y m\'aximo. \\
\hline
\end{longtable}

\begin{longtable}{|p{3cm}|p{11.5cm}|}
\hline
\multicolumn{2}{|c|}{\textbf{RNF-002} \criticomark} \\
\hline
Categor\'ia & Eficiencia de desempe\~no -- App m\'ovil \\
\hline
Obligatoriedad & \textbf{deber\'a} \\
\hline
Descripci\'on & Cat\'alogo de 16 especies carga en $\leq$3 segundos; navegaci\'on entre fichas en $\leq$0.5 segundos. \\
\hline
Fuente & \inferido{} \\
\hline
Prioridad & Alta \\
\hline
Verificaci\'on & Prueba de rendimiento con cat\'alogo de 16 especies. \\
\hline
\end{longtable}

\begin{longtable}{|p{3cm}|p{11.5cm}|}
\hline
\multicolumn{2}{|c|}{\textbf{RNF-021} \criticomark} \\
\hline
Categor\'ia & Eficiencia de desempe\~no -- Plataforma web \\
\hline
Obligatoriedad & \textbf{deber\'a} \\
\hline
Descripci\'on & La plataforma web deber\'a cumplir en conexi\'on de 10 Mbps:\newline
(a) Dashboard de validaci\'on ($\leq$50 observaciones): $\leq$3 segundos.\newline
(b) Mapa interactivo ($\leq$500 observaciones): $\leq$5 segundos.\newline
(c) Exportaci\'on de informe PDF ($\leq$200 observaciones): $\leq$15 segundos. \\
\hline
Fuente & \inferido{} -- Nielsen Norman Group ($\leq$10s para tareas complejas). \\
\hline
Prioridad & Alta \\
\hline
Verificaci\'on & Prueba con Lighthouse/WebPageTest en conexi\'on simulada de 10 Mbps. \\
\hline
\end{longtable}

\begin{longtable}{|p{3cm}|p{11.5cm}|}
\hline
\multicolumn{2}{|c|}{\textbf{RNF-003}} \\
\hline
Categor\'ia & Capacidad -- Backend servidor \\
\hline
Obligatoriedad & \textbf{deber\'ia} \\
\hline
Descripci\'on & El API Backend deber\'ia soportar: 40 dispositivos m\'oviles con sincronizaciones en lotes diferidos + hasta 5 usuarios web concurrentes (Validador/Bot\'anico/Admin), dentro del presupuesto de S/.500 para servidor. \\
\hline
Fuente & ANEXO 3. Ref. 10 (RNF-02). \\
\hline
Prioridad & Media \\
\hline
Verificaci\'on & Prueba de carga: 10 dispositivos sincronizando + 3 usuarios web concurrentes sin bloqueos en PostgreSQL. \\
\hline
\end{longtable}

\subsection{Usabilidad}

\begin{longtable}{|p{3cm}|p{11.5cm}|}
\hline
\multicolumn{2}{|c|}{\textbf{RNF-004} \criticomark} \\
\hline
Categor\'ia & Usabilidad -- Operabilidad -- App m\'ovil \\
\hline
Obligatoriedad & \textbf{deber\'a} \\
\hline
Descripci\'on & Interfaz operable con una sola mano; elementos interactivos del flujo principal en mitad inferior de la pantalla (zona del pulgar). \\
\hline
Fuente & \inferido{} \\
\hline
Prioridad & Alta \\
\hline
Verificaci\'on & Heur\'isticas Nielsen + prueba con 5 guardaparques. Criterio: \'exito $\geq$90\% con una mano. \\
\hline
\end{longtable}

\begin{longtable}{|p{3cm}|p{11.5cm}|}
\hline
\multicolumn{2}{|c|}{\textbf{RNF-005} \criticomark} \\
\hline
Categor\'ia & Usabilidad -- Aprendizaje -- App m\'ovil \\
\hline
Obligatoriedad & \textbf{deber\'a} \\
\hline
Descripci\'on & Flujo principal en m\'aximo 3 pasos (abrir $\rightarrow$ capturar $\rightarrow$ ver resultado). Flujo asistido: cada pregunta guiada es un paso adicional claramente numerado. \\
\hline
Fuente & CRISP-DM Fase 1. \\
\hline
Prioridad & Alta \\
\hline
Verificaci\'on & Task flow analysis + prueba primer uso sin entrenamiento. Criterio: identificaci\'on completa en $\leq$5 minutos. \\
\hline
\end{longtable}

\begin{longtable}{|p{3cm}|p{11.5cm}|}
\hline
\multicolumn{2}{|c|}{\textbf{RNF-006} \criticomark} \\
\hline
Categor\'ia & Usabilidad -- Accesibilidad -- App m\'ovil \\
\hline
Obligatoriedad & \textbf{deber\'a} \\
\hline
Descripci\'on & Botones t\'actiles $\geq$48$\times$48 dp; espaciado m\'inimo 8 dp entre elementos (uso con guantes de campo). \\
\hline
Fuente & \inferido{} -- Material Design de Google. \\
\hline
Prioridad & Alta \\
\hline
Verificaci\'on & Accessibility Scanner de Android. \\
\hline
\end{longtable}

\begin{longtable}{|p{3cm}|p{11.5cm}|}
\hline
\multicolumn{2}{|c|}{\textbf{RNF-007}} \\
\hline
Categor\'ia & Usabilidad -- Internacionalizaci\'on + Paleta corporativa \\
\hline
Obligatoriedad & \textbf{deber\'a} \\
\hline
Descripci\'on & Interfaz principal en espa\~nol. Nombres en quechua e ingl\'es con \'icono diferenciador. La plataforma web deber\'a implementar la \textbf{paleta de colores corporativa} definida por el equipo web (colores base: \texttt{\#757B47} y \texttt{\#FEFAE0} y tonos complementarios), superando puntuaci\'on de accesibilidad $>$90/100 en Google Lighthouse (Ref. 10, RNF-01). \\
\hline
Fuente & \inferido{}. Ref. 10 (RNF-01). \\
\hline
Prioridad & Media \\
\hline
Verificaci\'on & Audit\'ia Lighthouse en la plataforma web. Inspecci\'on de pantallas en app. \\
\hline
\end{longtable}

\begin{longtable}{|p{3cm}|p{11.5cm}|}
\hline
\multicolumn{2}{|c|}{\textbf{RNF-008}} \\
\hline
Categor\'ia & Usabilidad -- Accesibilidad visual -- App m\'ovil \\
\hline
Obligatoriedad & \textbf{deber\'ia} \\
\hline
Descripci\'on & Modo de alto contraste activable; contraste $\geq$4.5:1 (WCAG 2.1 AA) para uso bajo luz solar directa en el SHMP. \\
\hline
Fuente & \inferido{} \\
\hline
Prioridad & Media \\
\hline
Verificaci\'on & Prueba de contraste + prueba en campo con luz solar. \\
\hline
\end{longtable}

\subsection{Fiabilidad}

\begin{longtable}{|p{3cm}|p{11.5cm}|}
\hline
\multicolumn{2}{|c|}{\textbf{RNF-009} \criticomark} \\
\hline
Categor\'ia & Fiabilidad -- Disponibilidad -- App m\'ovil \\
\hline
Obligatoriedad & \textbf{deber\'a} \\
\hline
Descripci\'on & El 100\% de las funcionalidades de captura, identificaci\'on (incluyendo flujo asistido RF-049/050), cat\'alogo, mapa y registro deber\'an estar disponibles sin conexi\'on. \\
\hline
Fuente & CRISP-DM Fase 1. \\
\hline
Prioridad & Alta \\
\hline
Verificaci\'on & Prueba completa en modo avi\'on (20 identificaciones, incluyendo flujo asistido). \\
\hline
\end{longtable}

\begin{longtable}{|p{3cm}|p{11.5cm}|}
\hline
\multicolumn{2}{|c|}{\textbf{RNF-010} \criticomark} \\
\hline
Categor\'ia & Fiabilidad -- Recuperabilidad -- App m\'ovil \\
\hline
Obligatoriedad & \textbf{deber\'a} \\
\hline
Descripci\'on & Guardado autom\'atico de observaciones en progreso. Cero observaciones perdidas ante cierre inesperado, p\'erdida de bater\'ia o almacenamiento lleno. \\
\hline
Fuente & \inferido{} \\
\hline
Prioridad & Alta \\
\hline
Verificaci\'on & Inyecci\'on de fallos: cierre forzado durante 10 creaciones. Criterio: 0 observaciones perdidas. \\
\hline
\end{longtable}

\begin{longtable}{|p{3cm}|p{11.5cm}|}
\hline
\multicolumn{2}{|c|}{\textbf{RNF-011}} \\
\hline
Categor\'ia & Fiabilidad -- Tolerancia a fallos -- App m\'ovil + Backend \\
\hline
Obligatoriedad & \textbf{deber\'a} \\
\hline
Descripci\'on & Sincronizaci\'on con retroceso exponencial, m\'ax. 5 reintentos. Observaciones locales hasta confirmaci\'on del servidor. Descarga OTA del modelo soporta reanudaci\'on (RF-044). \\
\hline
Fuente & \inferido{} \\
\hline
Prioridad & Alta \\
\hline
Verificaci\'on & Prueba con desconexiones simuladas en sincronizaci\'on y descarga OTA. \\
\hline
\end{longtable}

\subsection{Seguridad}

\begin{longtable}{|p{3cm}|p{11.5cm}|}
\hline
\multicolumn{2}{|c|}{\textbf{RNF-012} \criticomark} \\
\hline
Categor\'ia & Seguridad -- Privacidad -- App m\'ovil + Web \\
\hline
Obligatoriedad & \textbf{deber\'a} \\
\hline
Descripci\'on & Sin recolecci\'on de datos sin consentimiento expl\'icito (Ley N.° 29733). Consentimiento en primer inicio (app) y primer acceso (web). \\
\hline
Fuente & ANEXO 2 (Componente \'Etico). \\
\hline
Prioridad & Alta \\
\hline
Verificaci\'on & Inspecci\'on de c\'odigo + revisi\'on legal. \\
\hline
\end{longtable}

\begin{longtable}{|p{3cm}|p{11.5cm}|}
\hline
\multicolumn{2}{|c|}{\textbf{RNF-013} \criticomark} \\
\hline
Categor\'ia & Seguridad -- Confidencialidad -- App m\'ovil \\
\hline
Obligatoriedad & \textbf{deber\'a} \\
\hline
Descripci\'on & Datos locales cifrados con Android Keystore + AES-256: observaciones, GPS, fotos y datos de sesi\'on. \\
\hline
Fuente & Ley N.° 29733. ANEXO 2. \\
\hline
Prioridad & Alta \\
\hline
Verificaci\'on & Extracci\'on BD local: verificar ilegibilidad sin clave Keystore. \\
\hline
\end{longtable}

\begin{longtable}{|p{3cm}|p{11.5cm}|}
\hline
\multicolumn{2}{|c|}{\textbf{RNF-014} \criticomark} \\
\hline
Categor\'ia & Seguridad -- Integridad -- App + Web + Backend \\
\hline
Obligatoriedad & \textbf{deber\'a} \\
\hline
Descripci\'on & Comunicaciones con el servidor via HTTPS/TLS. Contrase\~nas hasheadas con \textbf{Bcrypt} en PostgreSQL (factor $\geq$10) -- el hash es irreversible, no texto plano (Ref. 10, RNF-03). El API Backend deber\'a implementar pol\'iticas CORS y validar JWT en cada petici\'on protegida; sin token v\'alido: HTTP 401 (Ref. 10, RF-03). La plataforma web deber\'a implementar protecci\'on CSRF y cabeceras de seguridad (CSP, HSTS, X-Frame-Options). \\
\hline
Fuente & \inferido{}. Ref. 10 (RF-03, RNF-03). \\
\hline
Prioridad & Alta \\
\hline
Verificaci\'on & An\'alisis de tr\'afico + inspecci\'on BD (campo contrase\~na: hash Bcrypt, no texto plano) + prueba de petici\'on sin JWT (HTTP 401). \\
\hline
\end{longtable}

\subsection{Interoperabilidad}

\begin{longtable}{|p{3cm}|p{11.5cm}|}
\hline
\multicolumn{2}{|c|}{\textbf{RNF-022} \criticomark} \\
\hline
Categor\'ia & Interoperabilidad -- App m\'ovil $\leftrightarrow$ Web $\leftrightarrow$ Backend \\
\hline
Obligatoriedad & \textbf{deber\'a} \\
\hline
Descripci\'on & El API Backend (Node.js) deber\'a soportar lecturas/escrituras simult\'aneas de la plataforma web y de la app m\'ovil sin conflictos en PostgreSQL. Las peticiones concurrentes de ambos clientes no deber\'an generar bloqueos en la base de datos. \\
\hline
Fuente & Ref. 10 (RNF-02). \\
\hline
Prioridad & Alta \\
\hline
Verificaci\'on & Prueba de carga: peticiones concurrentes desde app m\'ovil y navegador web sin bloqueos en PostgreSQL. \\
\hline
\end{longtable}

\begin{longtable}{|p{3cm}|p{11.5cm}|}
\hline
\multicolumn{2}{|c|}{\textbf{RNF-023} \criticomark} \\
\hline
Categoría & Eficiencia de Memoria -- Optimización del Modelo IA \\
\hline
Obligatoriedad & \textbf{deberá} \\
\hline
Descripción & El modelo de IA (TensorFlow Lite) debe ser sometido a un proceso de cuantización de pesos (Int8) para reducir su tamaño y consumo de memoria RAM, garantizando que el proceso de inferencia no exceda los 150 MB de uso de memoria activa, permitiendo así la multitarea en dispositivos con 2 GB de RAM. \\
\hline
Fuente & Arquitectura Edge AI y restricciones de hardware. \\
\hline
Prioridad & Alta \\
\hline
Verificación & Perfilado de memoria durante 20 inferencias consecutivas; consumo pico ≤150 MB en dispositivo con 2 GB RAM. \\
\hline
\end{longtable}


\subsection{Portabilidad}

\begin{longtable}{|p{3cm}|p{11.5cm}|}
\hline
\multicolumn{2}{|c|}{\textbf{RNF-015} \criticomark} \\
\hline
Categor\'ia & Portabilidad -- Adaptabilidad -- App m\'ovil \\
\hline
Obligatoriedad & \textbf{deber\'a} \\
\hline
Descripci\'on & Compatible con Android 8.0 (API 26)+; 2 GB RAM; 16 GB almacenamiento. \\
\hline
Fuente & ANEXO 2. \\
\hline
Prioridad & Alta \\
\hline
Verificaci\'on & Prueba en Android 8.0, 10 y 13. \\
\hline
\end{longtable}

\begin{longtable}{|p{3cm}|p{11.5cm}|}
\hline
\multicolumn{2}{|c|}{\textbf{RNF-016}} \\
\hline
Categor\'ia & Portabilidad -- Instalabilidad -- App m\'ovil + Web \\
\hline
Obligatoriedad & \textbf{deber\'a} \\
\hline
Descripci\'on & App distribuida como APK sin Google Play Store. Plataforma web compatible con Chrome/Firefox/Edge (\'ultimas 2 versiones principales) sin instalaci\'on adicional. \\
\hline
Fuente & \inferido{}. Ref. 10. \\
\hline
Prioridad & Media \\
\hline
Verificaci\'on & Instalaci\'on APK sin cuenta Google. Prueba web en los 3 navegadores. \\
\hline
\end{longtable}

\subsection{Mantenibilidad}

\begin{longtable}{|p{3cm}|p{11.5cm}|}
\hline
\multicolumn{2}{|c|}{\textbf{RNF-017}} \\
\hline
Categor\'ia & Mantenibilidad -- Modularidad -- App + Web + Backend \\
\hline
Obligatoriedad & \textbf{deber\'ia} \\
\hline
Descripci\'on & App: capas separadas (UI, dominio, datos, IA). Web: componentes desacoplados (React, Node.js/Express, PostgreSQL). A\~nadir una especie al cat\'alogo (hasta 17+) no requiere modificar c\'odigo, solo datos. \\
\hline
Fuente & CRISP-DM Fase 6. Ref. 10. \\
\hline
Prioridad & Media \\
\hline
Verificaci\'on & Revisi\'on de arquitectura con diagrama de capas. \\
\hline
\end{longtable}

\subsection{Calidad del Modelo IA}

\begin{longtable}{|p{3cm}|p{11.5cm}|}
\hline
\multicolumn{2}{|c|}{\textbf{RNF-018} \criticomark} \\
\hline
Categor\'ia & Calidad del modelo CNN -- 16 especies \\
\hline
Obligatoriedad & \textbf{deber\'a} \\
\hline
Descripci\'on & El modelo CNN sobre las 16 especies deber\'a alcanzar:\newline
(a) \textbf{Accuracy global $\geq$85\%} sobre el dataset de prueba (15\% del total).\newline
(b) \textbf{Precision por especie $\geq$70\%}: ninguna especie con precision $<$70\%, cr\'itico para esp. CR/EN.\newline
(c) \textbf{Tasa de error en campo $<$10\%} en pruebas con guardaparques. \\
\hline
Fuente & ANEXO 2 (Hip\'otesis 1). CRISP-DM Fase 4. \\
\hline
Prioridad & Alta \\
\hline
Verificaci\'on & (a) Matriz de confusi\'on 16 clases. (b) Classification report con precision por clase. (c) 10 guardaparques $\times$ 30 espec\'imenes. \\
\hline
\end{longtable}

\begin{longtable}{|p{3cm}|p{11.5cm}|}
\hline
\multicolumn{2}{|c|}{\textbf{RNF-019} \textit{(Reubicado desde RF-016)}} \\
\hline
Categor\'ia & Calidad del modelo -- Cobertura fenol\'ogica \\
\hline
Obligatoriedad & \textbf{deber\'ia} \\
\hline
Descripci\'on & El modelo deber\'ia reconocer estados fenol\'ogicos (floraci\'on, crecimiento, marchito). Condicionado a $\geq$200 im\'agenes etiquetadas por estado por especie. \\
\hline
Fuente & ANEXO 2. CRISP-DM Fase 2. \\
\hline
Prioridad & Media (condicionada al dataset) \\
\hline
Verificaci\'on & Accuracy $\geq$70\% por estado si el dataset tiene suficientes muestras. \\
\hline
\end{longtable}

\subsection{Eficiencia Energ\'etica}

\begin{longtable}{|p{3cm}|p{11.5cm}|}
\hline
\multicolumn{2}{|c|}{\textbf{RNF-020} \criticomark} \\
\hline
Categor\'ia & Eficiencia Energ\'etica -- App m\'ovil \\
\hline
Obligatoriedad & \textbf{deber\'a} \\
\hline
Descripci\'on & 20 identificaciones en jornada de 8h: consumo $\leq$\textbf{15\% de la bater\'ia} del dispositivo de referencia (2 GB RAM, octa-core 1.8 GHz, bater\'ia 3,000 mAh). Estrategias:\newline
(a) GPS bajo demanda: activado solo en creaci\'on de observaci\'on, apagado tras obtener coordenadas o timeout.\newline
(b) C\'amara bajo demanda: visor activo solo durante el flujo de captura.\newline
(c) Inferencia CNN: modelo cargado una vez por sesi\'on; delegados NNAPI/GPU cuando disponible. \\
\hline
Fuente & CR-06. \\
\hline
Prioridad & Alta \\
\hline
Verificaci\'on & 20 identificaciones en dispositivo de referencia; medir con Android Battery Historian. Criterio: $\leq$15\% consumido. \\
\hline
\end{longtable}

%=============================================================
\section{Ap\'endices}
%=============================================================

\subsection{Ap\'endice A: Modelo de datos v5.2}

\textbf{Nota v5.2:} BD implementada en \textbf{PostgreSQL 15} en el servidor (accedida por Node.js/Express mediante Sequelize u ORM equivalente) y replicada localmente en \textbf{SQLite/Room} en la app m\'ovil. Cambios respecto a v5.1: tabla \texttt{Especie} con atributos visuales confirmados; tabla \texttt{ReporteEspecieDesconocida} y \texttt{RespuestaFlujGuiado} para el flujo asistido; RF-055 (cola del Bot\'anico); RF-052 (mapa web).

\begin{longtable}{|p{3.5cm}|p{11cm}|}
\hline
\textbf{Entidad} & \textbf{Atributos principales} \\
\hline
\endhead
\textbf{Usuario} & usuario\_id (PK), nombres, apellidos, email, password\_hash (Bcrypt), rol [guardaparque/validador/botanico/administrador], puesto\_control\_id (FK), telefono, activo, fecha\_registro, ultimo\_acceso, token\_local, fecha\_expiracion\_sesion \\
\hline
\textbf{PuestoControl} & puesto\_id (PK), nombre, zona\_shmp, latitud\_aprox, longitud\_aprox \\
\hline
\textbf{ZonaInteresSHMP} & zona\_id (PK), nombre [Aguas Calientes/Camino Inca/Zona Arqueol\'ogica/Bosque Montano/...], descripcion, tipo [sector/ruta/punto\_referencia], coordenadas\_geojson \\
\hline
\textbf{Especie} & especie\_id (PK), nombre\_cientifico, nombre\_comun\_es, nombre\_comun\_en, nombre\_quechua, biotipo (Ochoa 1), sistema\_ecologico (Ochoa 2), endemismo [local\_e*/nacional\_e] (Ochoa 3), estado\_conservacion\_iucn [CR/EN/VU/LC] (Ochoa 4), reporte\_nuevo [m/?] (Ochoa 5), descripcion\_es (Ochoa 7), detalle\_flor (Ochoa 8), descripcion\_en (Ochoa 11), habitat, epoca\_floracion, endemica\_peru (bool), color\_labelo, morfologia\_labelo [tubular/abierto/espatulado], tipo\_florecimiento [apical/lateral/basal/racemoso/solitario], morfologia\_foliar, disposicion\_hojas, biotipo\_detallado [epifita/terrestre/litofita], presencia\_camino\_inca (bool), presencia\_jardin\_botanico\_mp (bool), presencia\_jardin\_botanico\_llaqta (bool), validada (bool), validador\_id (FK), fecha\_validacion \\
\hline
\textbf{ImagenReferencia} & imagen\_id (PK), especie\_id (FK), ruta\_archivo, angulo [frontal\_labelo/lateral/hoja/general], estado\_fenologico, validador\_id (FK), fecha\_captura \\
\hline
\textbf{Observacion} & observacion\_id (PK), usuario\_id (FK), especie\_identificada\_id (FK nullable), fecha\_hora, latitud, longitud, precision\_gps\_m, puesto\_control\_id (FK), estado\_individuo [floracion/crecimiento/marchito], notas, confianza\_ia, estado\_sincronizacion, fecha\_sincronizacion, condicion\_climatica, microhabitat, datos\_campo\_completos (bool), intentos\_identificacion (int) \\
\hline
\textbf{ImagenObservacion} & imagen\_obs\_id (PK), observacion\_id (FK), ruta\_archivo, thumbnail\_ruta, angulo, es\_imagen\_principal (bool), fecha\_captura, foto\_purgada (bool) \\
\hline
\textbf{PreguntaGuiada} & pregunta\_id (PK), texto\_pregunta, opciones\_json, tipo [opciones/texto\_libre], activa (bool), orden \\
\hline
\textbf{RespuestaFlujGuiado} & respuesta\_id (PK), observacion\_id (FK), pregunta\_id (FK), respuesta\_opcion, respuesta\_libre\_texto, fecha\_respuesta \\
\hline
\textbf{ReporteEspecieDesconocida} & reporte\_id (PK), observacion\_id (FK), usuario\_id (FK), vector\_probabilidades\_json, estado [pendiente/asignada/nueva\_especie/descartada], especie\_asignada\_id (FK nullable), botanico\_revisor\_id (FK nullable), fecha\_revision, comentario\_botanico, alerta\_nueva\_especie (bool), fecha\_reporte \\
\hline
\textbf{Validacion} & validacion\_id (PK), observacion\_id (FK), validador\_id (FK), especie\_id\_corregida (FK nullable), estado\_validacion [aprobada/corregida/rechazada], comentario\_tecnico, justificacion\_rechazo, fecha\_validacion \\
\hline
\textbf{ModeloCNN} & modelo\_id (PK), version, ruta\_archivo\_tflite, checksum\_sha256, fecha\_publicacion, publicado\_por (FK), activo (bool), precision\_global, notas\_version, num\_especies (int) \\
\hline
\textbf{ConfiguracionRemota} & config\_id (PK), clave [umbral\_laplaciano/zoom\_cache\_max/umbral\_confianza\_ia/max\_intentos\_flujo\_asistido], valor, descripcion, fecha\_actualizacion \\
\hline
\textbf{Encuesta} & encuesta\_id (PK), usuario\_id (FK), fecha, precision\_percibida (1-5), utilidad\_percibida (1-5), facilidad\_uso (1-5), tiempo\_percibido (1-5), comentario\_libre \\
\hline
\end{longtable}

\textbf{Notas de implementaci\'on v5.2:}
\begin{itemize}
  \item \textbf{PostgreSQL (servidor):} contiene todas las entidades. Accedido por Node.js/Express mediante ORM (Sequelize recomendado). Las contrase\~nas se almacenan en \texttt{password\_hash} usando Bcrypt (factor $\geq$10) -- nunca texto plano.
  \item \textbf{SQLite/Room (app m\'ovil):} subconjunto local de \texttt{Usuario}, \texttt{Especie}, \texttt{Observacion}, \texttt{ImagenObservacion}, \texttt{PreguntaGuiada}, \texttt{RespuestaFlujGuiado}, \texttt{ReporteEspecieDesconocida} y \texttt{ConfiguracionRemota}. Cifrado AES-256 con SQLCipher + Android Keystore.
  \item \textbf{ZonaInteresSHMP:} nueva entidad para las capas de zonas del mapa (RF-035, RF-052). Contiene los 7 puestos de control, sectores del Santuario y rutas de patrullaje como geometr\'ias GeoJSON.
  \item \textbf{Rol en \texttt{Usuario}:} cuatro valores exactos -- \texttt{guardaparque}, \texttt{validador}, \texttt{botanico}, \texttt{administrador} -- alineados con los 4 actores del sistema.
\end{itemize}

\subsection{Ap\'endice B: Casos de uso principales}

\begin{longtable}{|p{2.3cm}|p{2.5cm}|p{2.8cm}|p{5.7cm}|}
\hline
\textbf{Actor} & \textbf{Plataforma} & \textbf{Acci\'on} & \textbf{Resultado esperado} \\
\hline
\endhead
Guardaparque & App m\'ovil & Identificar orqu\'idea & Resultado en $<$5s; top-3 con endemismo e IUCN destacados; sabe exactamente qu\'e est\'a viendo \\
\hline
Guardaparque & App m\'ovil & Baja confianza & Recomendaci\'on de nueva foto con \'angulo espec\'ifico (labelo, hoja) \\
\hline
Guardaparque & App m\'ovil & 2+ intentos fallidos & Preguntas guiadas: labelo, florecimiento, hojas, biotipo \\
\hline
Guardaparque & App m\'ovil & Reportar desconocida & Reporte con fotos, respuestas, GPS y datos de campo sincronizados \\
\hline
Guardaparque & App m\'ovil & Consultar cat\'alogo & Ficha con 11 atributos + labelo, florecimiento, hojas; funciona offline \\
\hline
Guardaparque & App m\'ovil & Ver mapa & Mapa con sus observaciones y zonas de inter\'es del SHMP \\
\hline
Validador & Web & Validar observaciones & Dashboard web con toda la informaci\'on biol\'ogica; aprobar/corregir/rechazar \\
\hline
Validador & Web & Ver mapa web & Mapa con todos los puntos, capas de puestos y zonas del Santuario \\
\hline
Bot\'anico & Web & Revisar reporte desconocida & Cola con fotos, respuestas y vector CNN; asignar / descartar / nueva especie \\
\hline
Bot\'anico & Web & Editar ficha especie & Completa 11 atributos Ochoa + atributos visuales (labelo, florecimiento, hojas) \\
\hline
Administrador & Web & Crear usuario & Formulario web $\rightarrow$ JSON $\rightarrow$ API (HTTP 201) $\rightarrow$ PostgreSQL \\
\hline
Administrador & Web & Actualizar modelo & Sube .tflite; distribuci\'on OTA con reanudaci\'on; verifica SHA-256 \\
\hline
Administrador & Web & Ver estad\'isticas & Panel global: usuarios, observaciones, reportes, versi\'on del modelo \\
\hline
\end{longtable}

\subsection{Ap\'endice C: Matriz de Trazabilidad RF -- Objetivos del Proyecto}

\begin{longtable}{|p{3cm}|p{5.2cm}|p{2.8cm}|p{2.5cm}|}
\hline
\textbf{RF/RNF} & \textbf{Objetivo / Restricci\'on} & \textbf{Fase CRISP-DM} & \textbf{Plataforma} \\
\hline
\endhead
RF-001 a RF-005 & RBAC; partici\'on web/app por rol; JWT; Bcrypt & Fase 6 & App+Web \\
\hline
RF-006 a RF-011 & OE4: calidad de imagen; preprocesamiento & Fases 3, 6 & App \\
\hline
RF-012 a RF-018 & OE3: CNN 16 especies; endemismo e IUCN visibles & Fase 4 & App \\
\hline
RF-049, RF-050 & Flujo asistido: labelo, florecimiento, hojas & Fases 4, 6 & App \\
\hline
RF-051 & Posible nueva especie; conservaci\'on & Fase 5 & App+Web \\
\hline
RF-019 a RF-023 & OE1+OE2: 16 fichas con atributos visuales & Fases 2, 3 & App+Web \\
\hline
RF-024 a RF-029 & OE4: registro biol\'ogico completo; purga; sync & Fase 6 & App \\
\hline
RF-030 a RF-034 & OE3: validaci\'on web; error 30\%$\to$10\% & Fase 5 & Web \\
\hline
RF-035, RF-036, RF-037, RF-038 & OE4: mapas con zonas de inter\'es; GPS degradado & Fase 6 & App+Web \\
\hline
RF-052 & OE4: mapa web con puestos y zonas del SHMP & Fase 6 & Web \\
\hline
RF-039 a RF-042 & OE5: informes PDF; m\'etricas; encuesta & Fase 5 & Web+App \\
\hline
RF-043 a RF-048 & Admin CRUD usuarios; OTA resumable; Ley 29733 & Fase 6 & Web \\
\hline
RF-055, RF-056 & Bot\'anico: cola especie desconocida; tendencias & Fases 2--5 & Web \\
\hline
RF-057, RF-058 & Compatibilidad API; JSON React$\leftrightarrow$Node.js & Fase 6 & Backend \\
\hline
RF-059 & Edge AI offline inmediato (<2s) & Fase 4 & App \\
\hline
RF-060 & Cola de sincronización asíncrona & Fase 6 & App+Backend \\
\hline
RNF-001, 002 & Hip\'otesis 2: $<$5 min identificaci\'on & Fase 1 & App \\
\hline
RNF-003 & Restricci\'on presupuestaria S/.500 & -- & Backend \\
\hline
RNF-004 a 008 & OE4: usabilidad campo; paleta corporativa web & Fase 1 & App+Web \\
\hline
RNF-009 a 011 & OE4: offline; tolerancia fallos; OTA & Fase 1 & App+Backend \\
\hline
RNF-012, 013, 014 & Ley 29733; Bcrypt; JWT; CORS & -- & App+Web+Backend \\
\hline
RNF-015, 016 & Android gama media; web multi-navegador & -- & App+Web \\
\hline
RNF-017 & Mantenibilidad; 17+ especies sin cambios c\'odigo & Fase 6 & App+Web \\
\hline
RNF-018 & OE3: accuracy 85\% + precision 70\%/especie & Fase 4 & Modelo IA \\
\hline
RNF-019 & OE1+OE2: cobertura fenol\'ogica & Fases 2--4 & Modelo IA \\
\hline
RNF-020 & Bater\'ia $\leq$15\% en 20 identificaciones & Fase 1 & App \\
\hline
RNF-021 & Tiempos de carga plataforma web & Fase 6 & Web \\
\hline
RNF-022 & Interoperabilidad App$\leftrightarrow$Web$\leftrightarrow$Backend & Fase 6 & Backend \\
\hline
RNF-023 & Cuantización Int8, memoria activa ≤150 MB & Fase 4 & App \\
\hline
\end{longtable}

\subsection{Ap\'endice D: Cronograma de desarrollo y RF por fase}

\begin{longtable}{|p{3.2cm}|p{3.8cm}|p{6.8cm}|}
\hline
\textbf{Etapa} & \textbf{Per\'iodo} & \textbf{RF/RNF relacionados} \\
\hline
Etapa 1: Literatura y dise\~no metodol\'ogico & Ago--Sep 2025 & RF-019 (16 fichas, atributos visuales); RF-050 (banco de preguntas guiadas) \\
\hline
Etapa 2: Recolecci\'on y procesamiento de datos & Oct 2025 -- Mar 2026 & RF-008/009; RF-023 (Bot\'anico valida 16 fichas); RNF-019 condicionado \\
\hline
Etapa 3: Desarrollo e integraci\'on & Mar--May 2026 & RF-001 a RF-058; RNF-001 a RNF-022. App m\'ovil + Plataforma web + API Backend \\
\hline
Etapa 4: Validaci\'on y cierre & May--Jul 2026 & RF-030--034 (validaci\'on); RF-052/055 (mapa web, cola Bot\'anico); RNF-018 (accuracy 16 especies); RNF-020 (bater\'ia); RNF-021 (web); RNF-022 (interoperabilidad) \\
\hline
\end{longtable}

\subsection{Ap\'endice E: Glosario del dominio}

\begin{longtable}{|p{3.5cm}|p{11cm}|}
\hline
\textbf{T\'ermino} & \textbf{Definici\'on} \\
\hline
\endhead
API Backend & Servidor Node.js 20 LTS + Express que act\'ua como n\'ucleo conector entre la app m\'ovil y la plataforma web. Expone endpoints REST JSON. \\
\hline
Bcrypt & Algoritmo de hashing de contrase\~nas utilizado en el API Backend (factor $\geq$10) para almacenar contrase\~nas de forma irreversible en PostgreSQL. \\
\hline
Biotipo & Forma de crecimiento de la orqu\'idea: ep\'ifita, terrestre, lit\'ofita. Atributo 1 Ochoa. \\
\hline
CORS & Cross-Origin Resource Sharing. Pol\'itica implementada en el API Backend para controlar qu\'e or\'igenes (app m\'ovil, plataforma web) pueden hacer peticiones. \\
\hline
Degradaci\'on Elegante & El sistema ofrece la mejor alternativa disponible sin bloquear al usuario (ej. GPS $>$10m). \\
\hline
Estado fenol\'ogico & Etapa del ciclo vital: floraci\'on, crecimiento vegetativo, marchito. \\
\hline
Ground truth & Datos validados por el Bot\'anico para entrenamiento y evaluaci\'on del modelo CNN. \\
\hline
JWT & JSON Web Token. Token firmado para autenticaci\'on entre clientes y el API Backend. \\
\hline
\textbf{Labelo} & P\'etalo modificado de las orqu\'ideas; elemento clave de identificaci\'on morfol\'ogica y visual. \\
\hline
Ochoa (11 atributos) & 11 campos de ficha de especie de la gu\'ia de Julio Gustavo Ochoa Estrada para el SHMP. \\
\hline
Paleta corporativa & Colores definidos por el equipo web para la plataforma: \texttt{\#757B47}, \texttt{\#FEFAE0} y tonos complementarios. \\
\hline
\textbf{Nueva Especie (posible)} & Planta no clasificable en las 16 especies tras m\'ultiples intentos y flujo asistido. \\
\hline
OTA & Over-The-Air. Actualizaci\'on del modelo CNN sin reinstalar la APK. \\
\hline
Resumabilidad & Capacidad de una descarga de reanudar desde el punto de interrupci\'on. \\
\hline
SERNANP & Servicio Nacional de \'Areas Naturales Protegidas por el Estado del Per\'u. \\
\hline
Sistema ecol\'ogico & Condiciones ecol\'ogicas del h\'abitat. Atributo 2 Ochoa. \\
\hline
TF Lite & TensorFlow Lite. Framework para inferencia de modelos ML en Android. \\
\hline
\textbf{Tipo de Florecimiento} & Caracter\'istica de la inflorescencia: apical, lateral, basal, racemoso, solitario. \\
\hline
Validador & Actor 2 del sistema (Especialista en RRNN). Accede solo a la plataforma web para validar observaciones. Distinto del Bot\'anico (Actor 3). \\
\hline
Zoom Level & Nivel de detalle del mapa. ZL 15/16 = sector/trocha del SHMP ($\sim$1:10,000--1:25,000). \\
\hline
Zona de inter\'es SHMP & Sector geogr\'afico del Santuario (Camino Inca, zona arqueol\'ogica, bosque montano, etc.) representado como capa en el mapa de la app y la web. \\
\hline
\end{longtable}

\subsection{Ap\'endice F: Registro de cambios v5.1 $\rightarrow$ v5.2}

\begin{longtable}{|p{1.5cm}|p{2.8cm}|p{2cm}|p{7.5cm}|}
\hline
\textbf{ID} & \textbf{Req. Afectado} & \textbf{Sev.} & \textbf{Descripci\'on} \\
\hline
\endhead
v5.2-01 & RF-005 & Alta & Roles renombrados y clarificados: Actor 2 pasa a denominarse ``Validador'' (antes Especialista) para distinguirlo del Bot\'anico (Actor 3). Descripci\'on de permisos actualizada. \\
\hline
v5.2-02 & Secci\'on 2.3 & Alta & Tablas de actores completamente reescritas con los 4 roles, plataforma expl\'icita y permisos RBAC precisos. El Administrador queda claramente como el \'unico con CRUD de usuarios. \\
\hline
v5.2-03 & RF-043 & Media & RF-043 actualizado: el Administrador gestiona cuentas de los 4 roles (guardaparque, validador, bot\'anico, administrador). Referencia expl\'icita a Bcrypt para contrase\~nas (Ref. 10, RNF-03). \\
\hline
v5.2-04 & RF-035 $\rightarrow$ RF-052 (nuevo) & Alta & Mapa de calor eliminado de v5.1. Reemplazado por: (1) RF-035 mejorado: mapa app con zonas de inter\'es del SHMP; (2) RF-052 nuevo: mapa interactivo web con capas de puestos de control, zonas y filtros. Sin an\'alisis de densidad ni detecci\'on de anomal\'ias. \\
\hline
v5.2-05 & RF-053, RF-054 (v5.1) & Alta & RF-053 (detecci\'on de eventos an\'omalos) y RF-054 (tendencias temporales) eliminados del SRS. Demasiado alcance para v1.0. RF-056 (tendencias temporales simplificado para Bot\'anico) conservado con prioridad Media. \\
\hline
v5.2-06 & RF-057, RF-058 (nuevos) & Alta & Dos nuevos RF de integraci\'on derivados del documento del equipo web (Ref. 10): RF-057 (compatibilidad API con app m\'ovil) y RF-058 (comunicaci\'on React$\leftrightarrow$Node.js en JSON). \\
\hline
v5.2-07 & RNF-022 (nuevo) & Alta & Nuevo RNF de interoperabilidad (Ref. 10, RNF-02): el API Backend soporta concurrencia de app m\'ovil y plataforma web sin bloqueos en PostgreSQL. \\
\hline
v5.2-08 & RNF-014 & Media & RNF-014 actualizado con requisitos de seguridad del API Backend: CORS, JWT (HTTP 401 sin token), Bcrypt (Ref. 10, RF-03, RNF-03). \\
\hline
v5.2-09 & RNF-007 & Media & RNF-007 actualizado con la paleta corporativa web (\texttt{\#757B47}, \texttt{\#FEFAE0}) y requisito de Lighthouse $>$90/100 (Ref. 10, RNF-01). \\
\hline
v5.2-10 & Ap\'endice A & Alta & Modelo de datos actualizado: rol \texttt{Usuario} con 4 valores exactos; tabla \texttt{ZonaInteresSHMP} a\~nadida para capas de mapa; nota de implementaci\'on PostgreSQL + Bcrypt para el servidor; Node.js/Express como ORM. \\
\hline
v5.2-11 & Stack tecnol\'ogico & Alta & Backend confirmado: Node.js 20 LTS + Express (no FastAPI). PostgreSQL 15 (servidor), SQLite/Room (app). React.js 18 + Tailwind CSS 3 (web). Axios/Fetch API como cliente HTTP. Alineado con Ref. 10. \\
\hline
\end{longtable}

\end{document}